\chapter{集成学习}
\label{chap:ensemble-learning}

在同一批信用风险数据上，两棵决策树可能因为训练样本稍有变化而形成不同的分支；一个线性模型与一棵树则可能分别抓住全局趋势和局部交互。面对这些差异，只从中选出一个模型，会丢掉其他模型已经学到的信息。能否让多个模型共同预测，并使组合在新样本上比单个成员更可靠？\term{集成学习}（ensemble learning）研究的正是成员怎样产生、怎样相互依赖，以及怎样把它们的输出组合起来。

组合并不自动带来收益。若成员反复犯相同的错误，增加数量只会重复同一判断；若成员彼此不同却普遍不准确，分歧本身也不能改善预测。有效的集成必须同时处理成员质量、错误依赖与组合规则。本章先建立这套共同分析框架，再分别讨论三类主要机制：Bagging通过训练扰动产生可并行的成员，Boosting让后续成员序贯修正当前组合，Stacking则从样本外预测中学习组合规则。

本章以监督预测为主要范围。组合后的回归响应、分类标签、得分和概率分别沿用第\ref{chap:regression}章与第\ref{chap:classification}章的任务语义；共识聚类见第\ref{sec:clustering-ensemble}节，弱学习与投票风险的证明见第\ref{sec:boosting}节。这里集中回答模型层的问题，不把成员数量、训练目标下降或某个软件实现的工程优势直接当作泛化保证。

\section{集成学习概述}
\label{sec:ensemble-overview}

\subsection{从单个模型到组合预测}
\label{sec:ensemble-from-single-to-combination}

第\ref{sec:regression-trees}节和第\ref{sec:classification-trees}节已经说明，树的分裂由有限训练样本决定。某些候选分裂的得分很接近时，替换少量样本就可能改变根节点，后续分支也随之变化。单棵树因而可能具有较大的训练波动。另一方面，一个受限的线性模型即使训练稳定，也可能无法表示风险只在局部条件组合下升高的现象。前一种困难主要来自估计不稳定，后一种困难主要来自表示受限；反复调节同一个模型的复杂度，不一定能同时消除二者。

集成学习在原有预测任务之外增加两个环节：先产生多个\term{成员模型}（member model），再组合这些成员的输出。强调训练方法或模型族时，也常把产生成员的算法称为\term{基学习器}（base learner）。成员可以是同一算法在不同样本或随机种子下的结果，也可以来自不同模型族；最终组合得到的对象称为\term{集成模型}（ensemble model）。这种做法没有改变响应或类别标签的含义，只改变了由训练数据形成预测器的路径。

组合方法围绕三个问题逐渐分化。Bagging从训练集重采样多个成员并固定聚合规则，使易受数据扰动影响的预测更稳定；Boosting让下一轮成员依赖当前组合，把新增容量分配给尚未处理好的训练信号；Stacking则用样本外成员预测训练另一层模型，使组合规则也可以学习。

图\ref{fig:ensemble-three-mechanisms}固定同一训练任务，比较三类方法改变的环节。Bagging中的成员可以彼此独立训练；Boosting的第$t$个成员依赖前$t-1$轮形成的当前预测；Stacking的基成员可以独立训练，但组合层必须看到不含当前样本训练信息的折外预测。这些差异比“用了多少个模型”更能说明训练成本与泄漏风险。

\begin{figure*}[htbp]
  \centering
  % Mechanism comparison for Bagging, Boosting, and Stacking.
% Schematic only; no empirical quantities are encoded.
% Canvas: 164 mm x 66 mm. Dependencies: project palette and TikZ.
\begin{tikzpicture}[
  x=1mm,y=1mm,
  font=\footnotesize,
  text=BookInk,
  source/.style={
    rectangle,
    draw=FigureInput,
    fill=FigureInput!8,
    line width=0.7pt,
    minimum width=29mm,
    minimum height=11mm,
    text width=25mm,
    align=center,
    inner sep=1.5mm
  },
  process/.style={
    rectangle,
    draw=FigureModel,
    fill=FigureModel!8,
    line width=0.7pt,
    minimum width=32mm,
    minimum height=11mm,
    text width=28mm,
    align=center,
    inner sep=1.5mm
  },
  combine/.style={
    rectangle,
    draw=FigureLoss,
    fill=FigureLoss!9,
    line width=0.8pt,
    minimum width=30mm,
    minimum height=11mm,
    text width=26mm,
    align=center,
    inner sep=1.5mm
  },
  output/.style={
    rectangle,
    draw=FigureOutput,
    fill=FigureOutput!6,
    line width=0.8pt,
    minimum width=18mm,
    minimum height=11mm,
    align=center,
    inner sep=1.5mm
  },
  flow/.style={
    -{Latex[length=2.2mm,width=1.5mm]},
    draw=BookMuted,
    line width=0.8pt,
    rounded corners=1.2mm
  },
  dependency/.style={
    -{Latex[length=2.2mm,width=1.5mm]},
    draw=FigureLoss,
    line width=0.8pt,
    dashed,
    rounded corners=1.2mm
  }
]
  \path[use as bounding box] (0,0) rectangle (164,66);

  \node[anchor=west,font=\heiti\small,text=BookInk] at (1,55) {Bagging};
  \node[source] (bag-source) at (40,55) {重采样或随机化};
  \node[process] (bag-members) at (82,55) {并行成员\\$h_1,\ldots,h_M$};
  \node[combine] (bag-combine) at (123,55) {固定平均或投票};
  \node[output] (bag-output) at (154,55) {$H$};
  \draw[flow] (bag-source) -- (bag-members);
  \draw[flow] (bag-members) -- (bag-combine);
  \draw[flow] (bag-combine) -- (bag-output);

  \node[anchor=west,font=\heiti\small,text=BookInk] at (1,33) {Boosting};
  \node[source] (boost-source) at (40,33) {当前组合$F_{t-1}$};
  \node[process] (boost-members) at (82,33) {序贯训练成员$h_t$};
  \node[combine] (boost-combine) at (123,33) {加法更新$F_t$};
  \node[output] (boost-output) at (154,33) {$F_T$};
  \draw[dependency] (boost-source) -- node[above=0.8mm,text=FigureLoss] {修正信号} (boost-members);
  \draw[flow] (boost-members) -- (boost-combine);
  \draw[flow] (boost-combine) -- (boost-output);
  \draw[dependency] (boost-combine.south) -- ++(0,-5)
    -| node[pos=0.35,above=0.8mm,text=FigureLoss] {下一轮} (boost-source.south);

  \node[anchor=west,font=\heiti\small,text=BookInk] at (1,11) {Stacking};
  \node[source] (stack-source) at (40,11) {折内训练基成员};
  \node[process] (stack-members) at (82,11) {折外预测矩阵$\bm{Z}$};
  \node[combine] (stack-combine) at (123,11) {训练元学习器$g$};
  \node[output] (stack-output) at (154,11) {$H$};
  \draw[flow] (stack-source) -- (stack-members);
  \draw[flow] (stack-members) -- (stack-combine);
  \draw[flow] (stack-combine) -- (stack-output);

  \draw[BookRule,line width=0.6pt] (0,44)--(164,44);
  \draw[BookRule,line width=0.6pt] (0,22)--(164,22);
\end{tikzpicture}

  \caption{Bagging、Boosting与Stacking的成员生成及组合关系。实线表示训练或预测的主路径，虚线表示由前一轮组合产生的修正信号。三类方法都输出组合预测，但成员间依赖与组合规则不同；图中流程为机制示意，不表示三者必须使用相同成员模型。}
  \label{fig:ensemble-three-mechanisms}
\end{figure*}

Bagging、Stacking和实用的AdaBoost分别由Breiman、Wolpert及Freund与Schapire的工作形成代表性方法，梯度提升又把序贯加法推广到一般损失\scite{breiman1996bagging,wolpert1992stacked,freund1997boosting,friedman2001gradientboost}。

\subsection{成员、组合规则与输出}
\label{sec:ensemble-members-combination-output}

设训练集为
\[
  D_n=\{(\bm{x}_i,y_i)\}_{i=1}^{n},
\]
学习过程产生$M$个成员模型$h_1,\ldots,h_M$。用$\mathcal{C}$表示组合规则，集成预测统一写成
\begin{equation}
  H(\bm{x})
  =\mathcal{C}\bigl(h_1(\bm{x}),\ldots,h_M(\bm{x})\bigr).
  \label{eq:ensemble-combination}
\end{equation}
式\eqref{eq:ensemble-combination}只规定组合接口，还没有规定输出的语义。回归成员若都预测同一个数值目标，可以采用等权平均
\[
  H(\bm{x})=\frac{1}{M}\sum_{m=1}^{M}h_m(\bm{x}),
\]
也可以用非负权重$w_m$作加权平均，其中$\sum_m w_m=1$。分类成员则可能输出硬标签、类别得分或类别概率。硬标签适合投票，得分和概率可以逐类平均后再决策，但三种量不能在没有转换的情况下混合。

输出类型之所以重要，是因为相同的代数操作可能对应不同结论。两个分类器的得分尺度不同，直接平均会让数值范围较大的成员取得更大影响；两个成员都输出概率，也不意味着平均后已经校准。若元学习器把成员得分映射为概率，还要单独检查这个映射在新样本上的可靠性。第\ref{sec:classification-evaluation}节给出的硬标签、排序得分和类别概率评价接口同样适用于集成模型。

组合权重也不是成员的固有属性。等权平均预先固定权重，AdaBoost在序贯训练中同时确定成员贡献，线性Stacking则根据折外预测估计权重。权重还可以依赖输入，此时$w_m=w_m(\bm{x})$，不同样本会调用不同成员或采用不同组合比例。成员数量、权重约束、组合层复杂度以及是否允许输入相关组合，都属于训练或验证流程需要确定的模型选择。

\subsection{成员质量、错误依赖与多样性}
\label{sec:ensemble-quality-dependence-diversity}

平均为什么可能比单个成员稳定，可以先看一个受控情形。固定输入$\bm{x}$，把训练集抽取、重采样和训练随机性都视为随机来源。若$M$个成员预测具有相同方差$\sigma^2(\bm{x})$，任意两个不同成员的预测相关系数均为$\rho(\bm{x})$，则等权平均的方差为
\begin{equation}
  \operatorname{Var}\!\left[H(\bm{x})\right]
  =
  \sigma^2(\bm{x})
  \left[
    \rho(\bm{x})
    +
    \frac{1-\rho(\bm{x})}{M}
  \right].
  \label{eq:ensemble-correlated-variance}
\end{equation}
其中随$M$衰减的是成员不共享的波动；$\rho(\bm{x})\sigma^2(\bm{x})$所代表的共同变化不会因继续增加同类成员而消失。式\eqref{eq:ensemble-correlated-variance}依赖等方差、等相关的简化设定，用来解释相关性的作用，不是任意集成的精确风险公式。

图\ref{fig:ensemble-error-dependence}把这一差别写成成员误差的两部分。成员特有波动$u_m$在平均中可能相互抵消；若每个成员都带有共同误差$c$，增加同类成员只能继续削弱前一部分，不能消除$c$。因此，成员数增长与成员错误依赖下降是两个不同的条件。

\begin{figure}[htbp]
  \centering
  % Averaging suppresses member-specific fluctuations but preserves shared error.
% Schematic only; no empirical quantities are encoded.
% Canvas: 112 mm x 57 mm. Dependencies: project palette and TikZ.
\begin{tikzpicture}[
  x=1mm,y=1mm,
  font=\footnotesize,
  text=BookInk,
  member/.style={
    circle,
    draw=FigureModel,
    fill=FigureModel!8,
    line width=0.7pt,
    minimum size=8mm,
    inner sep=0.5mm
  },
  shared/.style={
    rectangle,
    draw=FigureLoss,
    fill=FigureLoss!9,
    line width=0.8pt,
    minimum width=42mm,
    minimum height=8mm,
    align=center,
    inner sep=1mm
  },
  result/.style={
    rectangle,
    draw=FigureOutput,
    fill=FigureOutput!6,
    line width=0.8pt,
    minimum width=42mm,
    minimum height=10mm,
    text width=38mm,
    align=center,
    inner sep=1mm
  },
  flow/.style={
    -{Latex[length=2.2mm,width=1.5mm]},
    draw=BookMuted,
    line width=0.8pt,
    rounded corners=1.2mm
  }
]
  \path[use as bounding box] (0,0) rectangle (112,57);

  \node[anchor=west,font=\heiti\small] at (1,53)
    {(a) 只有成员特有波动};
  \foreach \x/\name/\lab in {7/u1/$u_1$,17/u2/$u_2$,27/u3/$u_3$,37/u4/$u_4$,47/u5/$u_M$} {
    \node[member] (\name) at (\x,38) {\lab};
  }
  \node[text=BookMuted,inner sep=0pt] at (42,38) {$\cdots$};
  \node[result] (independent-average) at (27,14)
    {平均后只剩\\$\displaystyle M^{-1}\sum_m u_m$};
  \foreach \name in {u1,u2,u3,u4,u5} {
    \draw[flow] (\name.south) -- ++(0,-4) -| (independent-average.north);
  }
  \node[anchor=north,text=FigureModel,align=center] at (27,7)
    {$M$增大时，不共享的波动可被削弱};

  \draw[BookRule,line width=0.6pt] (56,3)--(56,54);

  \node[anchor=west,font=\heiti\small] at (60,53)
    {(b) 成员还共享同一偏移};
  \node[shared] (common) at (86,42)
    {共同误差$c$进入每个成员};
  \node[member] (cu1) at (64,29) {$u_1$};
  \node[member] (cu2) at (75,29) {$u_2$};
  \node[text=BookMuted,inner sep=0pt] at (86,29) {$\cdots$};
  \node[member] (cu4) at (97,29) {$u_{M-1}$};
  \node[member] (cu5) at (108,29) {$u_M$};
  \node[result] (correlated-average) at (86,11)
    {平均误差\\$\displaystyle c+M^{-1}\sum_m u_m$};
  \draw[flow,draw=FigureLoss] (common.east) -- (111,42)
    |- (correlated-average.east);
  \foreach \name in {cu1,cu2,cu4,cu5} {
    \draw[flow] (\name.south) -- ++(0,-3) -| (correlated-average.north);
  }
  \node[anchor=north,text=FigureLoss,align=center] at (86,4)
    {共同误差$c$不会因增加同类成员而消失};
\end{tikzpicture}

  \caption{成员误差依赖与平均的作用。左图表示只有成员特有波动时，平均可以随成员数增加而削弱这些波动；右图表示所有成员还共享误差$c$时，共同项仍保留在平均预测中。图中成员数量与符号只用于说明机制，不表示实测误差。}
  \label{fig:ensemble-error-dependence}
\end{figure}

平方误差还给出另一个逐点恒等式。对非负权重$w_m$且$\sum_mw_m=1$，令$H(\bm{x})=\sum_mw_mh_m(\bm{x})$，则
\begin{equation}
  \bigl(y-H(\bm{x})\bigr)^2
  =
  \sum_{m=1}^{M}w_m\bigl(y-h_m(\bm{x})\bigr)^2
  -
  \sum_{m=1}^{M}w_m\bigl(h_m(\bm{x})-H(\bm{x})\bigr)^2.
  \label{eq:ensemble-ambiguity-decomposition}
\end{equation}
右侧第一项是成员平方误差的加权平均，第二项衡量成员围绕组合预测的分散程度，常称为歧义项\scite{krogh1995ensembles}。在成员平均误差固定时，更大的有效分散会降低组合误差；但若为了制造分散而加入很差的成员，第一项也会增大。多样性只能与成员质量一起解释。

分类投票没有完全相同的平方误差分解，却面对同一问题：关键不只是每个成员错多少次，还包括它们是否在同一批样本上出错。对一对分类器，可以记录二者同时正确、只有一个正确和同时错误的样本数$n_{1,1},n_{1,0},n_{0,1},n_{0,0}$。分歧率$(n_{1,0}+n_{0,1})/n$描述二者给出不同正确性结果的频率，双错率$n_{0,0}/n$直接关注共同错误；Q统计量
\[
  Q=
  \frac{n_{1,1}n_{0,0}-n_{1,0}n_{0,1}}
       {n_{1,1}n_{0,0}+n_{1,0}n_{0,1}}
\]
在分母非零时描述两种正确性结果的关联。不同指标关注的对象不同，也不存在一个脱离成员准确率、对所有任务都能单调预测集成风险的多样性分数\scite{kuncheva2003diversity}。

成员差异可以从多个位置产生。改变训练样本会改变成员看到的经验分布；改变输入特征会改变可用信息；改变模型结构或超参数会改变表示与正则；改变初始化、样本顺序或分裂候选会改变优化路径。Bagging主要扰动训练样本，随机森林和Extra Trees进一步扰动特征或阈值，Stacking则允许直接组合不同模型族。这些机制只是多样性的来源，最终是否形成互补仍要用样本外预测检查。

由此可以用三个维度比较后续方法：成员怎样产生，成员之间有什么依赖，组合规则怎样确定。Bagging采用扰动后的成员和固定聚合，Boosting采用序贯依赖成员和加法组合，Stacking采用样本外成员预测和学习式组合。三类方法都可能改善预测，也都可能因共同失配、数据泄漏或过度选择而失效。

\section{基于Bagging的集成}
\label{sec:ensemble-bagging}

若一个学习算法对训练样本的细小变化很敏感，可以在同一总体上设想反复取得训练集、反复训练模型，再平均预测。真实任务通常只有一份训练集，无法直接完成这种重复抽样。\term{Bagging}是bootstrap aggregating的缩写：它用自助采样从现有训练集构造多份样本，分别训练成员，再用固定规则聚合。该方法主要减弱能够被重采样暴露出来的训练波动，不会自动修正所有成员共享的表示偏差\scite{breiman1996bagging}。

\subsection{Bagging}
\label{sec:ensemble-bagging-bootstrap}

\subsubsection{自助样本与聚合预测}

对第$m$个成员，从$\{1,\ldots,n\}$中独立、等概率、有放回地抽取$n$次索引，所得多重集合记为$I^{(m)}$。同一原始样本可能出现多次，也可能一次都没有出现。学习算法$\mathcal{A}$在相应自助样本上训练出
\[
  h_m=\mathcal{A}\!\left(D_n^{(m)}\right),
  \qquad
  D_n^{(m)}=\{(\bm{x}_i,y_i):i\in I^{(m)}\}.
\]
回归通常平均$h_m(\bm{x})$；分类可以投票，也可以在成员输出具有同一语义时平均逐类得分或概率。自助样本大小不必固定为$n$，成员算法也不必是树，但方法是否受益取决于重采样能否产生有用的预测差异。

一份大小为$n$的标准自助样本没有抽中原始样本$i$的概率为
\begin{equation}
  \mathbb{P}\!\left(i\notin I^{(m)}\right)
  =
  \left(1-\frac{1}{n}\right)^n
  \longrightarrow e^{-1}.
  \label{eq:ensemble-oob-probability}
\end{equation}
因此，每个成员平均约有$36.8\%$的原始样本未参与其训练。这些样本构成该成员的\term{袋外样本}（out-of-bag sample, OOB sample）。反过来，对每个训练样本$i$，可以只收集没有使用它训练的成员预测。令
\[
  \mathcal{M}_{i}^{\mathrm{OOB}}
  =\{m:i\notin I^{(m)}\},
\]
当$\mathcal{M}_{i}^{\mathrm{OOB}}$非空时，回归的袋外预测可写为
\begin{equation}
  \widehat{y}_{i}^{\mathrm{OOB}}
  =
  \frac{1}{|\mathcal{M}_{i}^{\mathrm{OOB}}|}
  \sum_{m\in\mathcal{M}_{i}^{\mathrm{OOB}}}h_m(\bm{x}_i),
  \label{eq:ensemble-oob-prediction}
\end{equation}
分类则对同一成员集合投票或平均类别输出。图\ref{fig:ensemble-bagging-oob}强调了袋外预测的样本级条件：不同训练样本由不同成员子集预测，而不是另行训练一个统一的“OOB模型”。

\begin{figure*}[htbp]
  \centering
  % Bagging training and sample-specific out-of-bag prediction.
% Schematic only; no empirical quantities are encoded.
% Canvas: 164 mm x 66 mm. Dependencies: project palette and TikZ.
\begin{tikzpicture}[
  x=1mm,y=1mm,
  font=\figurefont\footnotesize,
  text=BookInk,
  data/.style={
    rectangle,
    model input,
    minimum width=20mm,
    minimum height=9mm,
    align=center,
    inner sep=1.2mm
  },
  model/.style={
    rectangle,
    model hidden,
    minimum width=16mm,
    minimum height=9mm,
    align=center,
    inner sep=1.2mm
  },
  filter/.style={
    rectangle,
    model training,
    minimum width=26mm,
    minimum height=12mm,
    text width=22mm,
    align=center,
    inner sep=1.5mm
  },
  result/.style={
    rectangle,
    model output,
    minimum width=22mm,
    minimum height=12mm,
    text width=21mm,
    align=center,
    inner sep=1.5mm
  },
  flow/.style={
    model flow
  },
  oob/.style={
    model feedback
  }
]
  \path[use as bounding box] (0,0) rectangle (164,66);

  \node[anchor=west,font=\figurefont\bfseries\small] at (1,62) {(a) 训练成员};
  \node[data] (dataset) at (12,32) {训练集\\$D_n$};

  \node[data] (sample1) at (39,51) {$D_n^{(1)}$};
  \node[data] (sample2) at (39,32) {$D_n^{(2)}$};
  \node[data] (samplem) at (39,13) {$D_n^{(M)}$};
  \node[text=BookMuted,inner sep=0pt] at (39,22.5) {$\vdots$};

  \node[model] (model1) at (67,51) {$h_1$};
  \node[model] (model2) at (67,32) {$h_2$};
  \node[model] (modelm) at (67,13) {$h_M$};
  \node[text=BookMuted,inner sep=0pt] at (67,22.5) {$\vdots$};

  \draw[flow] (dataset.east) -- ++(4,0) |- (sample1.west);
  \draw[flow] (dataset.east) -- (sample2.west);
  \draw[flow] (dataset.east) -- ++(4,0) |- (samplem.west);
  \draw[flow] (sample1) -- (model1);
  \draw[flow] (sample2) -- (model2);
  \draw[flow] (samplem) -- (modelm);
  \node[anchor=south,text=NNInput,inner sep=0pt] at (39,57)
    {有放回抽样};

  \draw[BookRule,line width=0.6pt] (81,4)--(81,62);

  \node[anchor=west,font=\figurefont\bfseries\small] at (86,62) {(b) 样本$i$的袋外预测};
  \node[data,minimum width=16mm] (query) at (94,33) {样本\\$\bm{x}_i$};
  \node[filter] (select) at (121,33)
    {筛选未使用$i$训练的成员\\
     $m\in\mathcal{M}_{i}^{\mathrm{OOB}}$};
  \node[result] (aggregate) at (151,33)
    {聚合预测\\$\widehat{y}_i^{\mathrm{OOB}}$};

  \draw[oob] (query) -- (select);
  \draw[flow] (select) -- (aggregate);
  \node[anchor=north,text=BookMuted,align=center,inner sep=0pt]
    at (121,22)
    {成员集合随样本$i$改变；\\有限$M$下集合大小也会变化};
\end{tikzpicture}

  \caption{Bagging训练与袋外预测。每个成员只使用其自助样本训练；对训练样本$\bm{x}_i$，袋外预测只聚合未抽中$i$的成员。图中索引和成员集合仅用于说明数据流；有限成员下，各样本可用的袋外成员数并不相同。}
  \label{fig:ensemble-bagging-oob}
\end{figure*}

\subsubsection{成员数量与误差来源}

Bagging的成员可以独立训练，因而适合把成员分配到不同计算单元。并行性只缩短墙钟时间，不会消除总计算量：若单个成员的训练和预测成本分别为$C_{\mathrm{train}}$与$C_{\mathrm{pred}}$，朴素实现的总工作量仍约随$M C_{\mathrm{train}}$和$M C_{\mathrm{pred}}$增长。模型存储也通常随成员数量增加。

增加成员数主要减小有限次随机抽样造成的Monte Carlo波动。若成员由同一随机训练机制产生，等权平均会随$M$增加逼近该随机机制下的期望预测；它不会消除原训练集有限、模型族不合适或成员高度相关造成的误差。实践中可以观察袋外损失随成员数的变化，在曲线趋于稳定后停止增加成员，但这个停止规则也属于模型选择过程。

成员模型的复杂度与成员数量承担不同职责。较深的树通常更能拟合每份自助样本，也更容易因样本扰动产生不同结构；平均可以削弱部分波动，却不意味着每棵树都应无条件生长到纯叶。叶最小样本量、最大深度、类别权重和自助样本大小仍须与成员数共同选择，尤其在标签噪声、极小叶或类别极不平衡时。

\subsubsection{袋外预测的用途与边界}

袋外预测没有使用当前样本训练相应成员，因而可用于观察成员数量、估计内部误差或构造置换重要性。它省去了为每个成员另设验证集的开销，但与$K$折交叉验证并不相同：每个样本的袋外成员数随机变化，成员训练集之间也高度重叠。成员数很小时，某些样本可能没有足够的袋外预测。

如果模型结构、成员数和特征处理都反复根据袋外结果选择，最终袋外误差已经参与了选择，不再是冻结流程后的独立测试结果。时间序列、分组样本或空间相关数据还要求自助抽样尊重依赖单位；逐行重采样会破坏原有数据结构。袋外机制提供的是一种内部样本外预测来源，而不是脱离数据生成条件的无偏泛化保证。

\subsection{随机森林}
\label{sec:ensemble-random-forest}

Bagging若以树为成员，各棵树仍可能反复选择少数强特征，产生相近的上层分裂。\term{随机森林}（random forest）在树的自助采样之外加入节点级特征随机化：每个节点先随机抽取一组候选特征，再只在这组特征中搜索分裂。树仍按第\ref{sec:regression-trees}节或第\ref{sec:classification-trees}节的任务准则形成叶输出，改变的是每个成员可参与分裂搜索的信息\scite{breiman2001randomforests}。

候选特征数控制单树质量与树间依赖的取舍。候选数接近全部特征时，森林更接近普通的树Bagging；候选数较小时，不同树更可能采用不同分裂，但单棵树也更容易错过当前节点上的有效特征。这一取舍还会受到特征冗余、样本量和树深影响，不能由一个固定默认值替代验证。

树数增加时，随机森林预测会逐渐稳定到其随机化机制对应的有限样本极限。这种稳定是对抽样随机性的收敛，不表示模型已经得到真实风险最小解，也不表示继续加树没有资源代价。随机森林的一致性结论通常针对特定的树构造、抽样比例与分布条件；实际软件在缺失值、类别特征、停止规则和并行实现上的差异，可能不满足同一组理论假设\scite{scornet2015consistency}。

\subsubsection{特征重要性的解释}

树集成常附带两类特征重要性。\term{不纯度下降重要性}（impurity-decrease importance）累计某个特征参与分裂时带来的加权目标下降，计算方便，却可能偏向候选切分较多的特征。它还会把相关特征之间的可替代关系压缩进实际被选中的分裂，不能解释为变量的独立贡献。

\term{置换重要性}（permutation importance）固定已训练模型，在评价数据上打乱一个特征，再观察预测损失增加多少。这个量直接对应已拟合模型在所选扰动下对该特征的依赖，但打乱会破坏该特征与其他特征的联合关系。高度相关的特征可能互相替代，使单个特征的置换影响看起来很小；不现实的特征组合也可能夸大损失变化。两类重要性都不是因果效应，需要连同评价数据、损失和置换方式一起报告\scite{strobl2007bias}。

图\ref{fig:ensemble-tree-randomization}沿一棵树的构造路径比较三种随机化位置。树Bagging主要改变成员看到的训练样本；随机森林还在每个节点抽取候选特征；Extra Trees继续把随机性放入候选阈值。三者都保留树的任务准则，但交给贪心搜索的信息逐步受到更多约束。

\begin{figure*}[htbp]
  \centering
  % Randomization locations in tree Bagging, random forests, and Extra Trees.
% Schematic only; no empirical quantities are encoded.
% Canvas: 164 mm x 69 mm. Dependencies: project palette and TikZ.
\begin{tikzpicture}[
  x=1mm,y=1mm,
  font=\footnotesize,
  text=BookInk,
  stage/.style={
    rectangle,
    draw=BookMuted,
    fill=BookPaper,
    line width=0.6pt,
    minimum width=39mm,
    minimum height=9mm,
    text width=35mm,
    align=center,
    inner sep=1mm
  },
  randomized/.style={
    stage,
    draw=FigureInput,
    fill=FigureInput!10,
    line width=0.9pt
  },
  searched/.style={
    stage,
    draw=FigureModel,
    fill=FigureModel!8,
    line width=0.8pt
  },
  output/.style={
    stage,
    draw=FigureOutput,
    fill=FigureOutput!6,
    line width=0.8pt
  },
  flow/.style={
    -{Latex[length=2.2mm,width=1.5mm]},
    draw=BookMuted,
    line width=0.8pt
  }
]
  \path[use as bounding box] (0,0) rectangle (164,69);

  \node[font=\heiti\small] at (27,65) {(a) 树Bagging};
  \node[randomized] (bag-sample) at (27,51) {随机化训练样本\\自助采样};
  \node[searched] (bag-feature) at (27,35) {检查全部候选特征};
  \node[searched] (bag-threshold) at (27,19) {搜索候选阈值};
  \node[output] (bag-tree) at (27,5) {得到一棵树};
  \draw[flow] (bag-sample)--(bag-feature);
  \draw[flow] (bag-feature)--(bag-threshold);
  \draw[flow] (bag-threshold)--(bag-tree);

  \node[font=\heiti\small] at (82,65) {(b) 随机森林};
  \node[randomized] (rf-sample) at (82,51) {随机化训练样本\\通常为自助采样};
  \node[randomized] (rf-feature) at (82,35) {随机化候选特征\\每个节点重新抽取};
  \node[searched] (rf-threshold) at (82,19) {在特征子集内搜索阈值};
  \node[output] (rf-tree) at (82,5) {得到一棵树};
  \draw[flow] (rf-sample)--(rf-feature);
  \draw[flow] (rf-feature)--(rf-threshold);
  \draw[flow] (rf-threshold)--(rf-tree);

  \node[font=\heiti\small] at (137,65) {(c) Extra Trees};
  \node[stage] (extra-sample) at (137,51)
    {训练样本\\原始方法使用完整训练集};
  \node[randomized] (extra-feature) at (137,35)
    {随机化候选特征};
  \node[randomized] (extra-threshold) at (137,19)
    {随机生成候选阈值\\再按准则择优};
  \node[output] (extra-tree) at (137,5) {得到一棵树};
  \draw[flow] (extra-sample)--(extra-feature);
  \draw[flow] (extra-feature)--(extra-threshold);
  \draw[flow] (extra-threshold)--(extra-tree);

  \draw[BookRule,line width=0.6pt] (54.5,1)--(54.5,66);
  \draw[BookRule,line width=0.6pt] (109.5,1)--(109.5,66);
\end{tikzpicture}

  \caption{树Bagging、随机森林与Extra Trees的随机化位置。蓝色节点表示引入随机性的环节，青色节点表示仍按任务准则执行的搜索，深色节点表示得到的成员树。Extra Trees面板采用原始方法的完整训练集设定；具体实现也可能额外启用自助采样。}
  \label{fig:ensemble-tree-randomization}
\end{figure*}

\subsection{极度随机树（Extra Trees）}
\label{sec:ensemble-extra-trees}

随机森林仍会在候选特征内搜索较优阈值。\term{极度随机树}（extremely randomized trees, Extra Trees）把随机性进一步放入阈值选择：在每个节点，对随机选出的候选特征生成随机切分点，再从这些随机候选中选择得分较好的切分。原始方法通常让各棵树使用完整训练集，把成员差异主要交给特征与阈值随机化；具体软件也可能允许额外使用自助采样\scite{geurts2006extra}。

随机阈值减少了穷举候选切分的工作，也会让单棵树偏离当前样本上的最佳贪心分裂。组合是否改善，取决于单树预测能力下降与成员依赖减弱之间的平衡。候选特征数、每个特征的随机阈值数、叶最小样本量以及是否采用自助采样共同决定这个平衡；“随机性更强”不能推出测试风险必然更低。

Extra Trees还说明，Bagging家族的共同点不是必须采用标准自助采样，而是通过受控扰动产生一组成员，并用固定规则聚合。随机子空间只扰动成员可见的特征，Pasting以无放回子样本替代自助样本，随机权重则让同一记录在不同成员中具有不同影响。每种扩展都应明确扰动发生在哪里，以及它是否保留袋外预测等附带机制。

\subsection{模型对比与总结}
\label{sec:ensemble-bagging-comparison}

表\ref{tab:ensemble-bagging-comparison}固定树为成员，比较三种方法怎样产生差异。普通Bagging只改变训练样本；随机森林再限制节点可见的候选特征；Extra Trees进一步随机化候选阈值。后两者并非简单地沿一条“随机性越多越好”的顺序排列，因为它们还可能采用不同抽样方案，并改变单树质量和搜索成本。

\begin{table*}[htbp]
  \centering
  \small
  \begin{tabularx}{\linewidth}{@{}>{\raggedright\arraybackslash}p{24mm}*{6}{>{\raggedright\arraybackslash}X}@{}}
    \toprule
    方法 & 样本扰动 & 节点候选特征 & 阈值选择 & 袋外预测 & 成员训练 & 主要边界 \\
    \midrule
    树Bagging & 通常自助采样 & 全部特征 & 在候选中搜索 & 标准自助设定下可用 & 成员间可并行 & 强特征可能使树高度相关 \\
    随机森林 & 通常自助采样 & 随机特征子集 & 在子集内搜索 & 标准自助设定下可用 & 成员间可并行 & 候选数过小会削弱单树 \\
    Extra Trees & 原始方法使用完整训练集；实现可选自助采样 & 随机特征子集 & 生成随机阈值后择优 & 仅在采用自助采样时可用 & 成员间可并行 & 过强随机化会增大单树误差 \\
    \bottomrule
  \end{tabularx}
  \caption{以树为成员时，Bagging、随机森林与Extra Trees的机制比较。表中描述稳定的方法定义；软件默认值可能随实现和版本改变，实际比较还须固定树复杂度、成员数与计算预算。}
  \label{tab:ensemble-bagging-comparison}
\end{table*}

Bagging适合把训练扰动转化为成员差异，再通过聚合降低其中不共享的部分。若所有成员都遗漏同一结构，固定聚合不会主动告诉下一棵树应修正什么。下一节转向序贯成员：新成员不再独立产生，而是根据当前组合留下的训练信号继续更新。

\section{基于Boosting的集成}
\label{sec:ensemble-boosting}

Bagging中的成员可以同时训练，因为一个成员不需要知道其他成员的预测。若希望新增成员专门补充当前组合尚未处理好的部分，这种独立性就不再成立。\term{Boosting}（提升）构造一个序贯加法模型：第$t$轮先观察已有函数$F_{t-1}$，再选择新成员及其贡献，使指定训练目标继续下降。AdaBoost、梯度提升及现代提升树共享这种依赖关系，但它们生成修正信号和近似目标的方式不同。

\subsection{前向逐步加法模型}
\label{sec:ensemble-stagewise-additive}

设$F_0$为初始预测，Boosting经过$T$轮得到
\begin{equation}
  F_T(\bm{x})
  =
  F_0(\bm{x})
  +
  \sum_{t=1}^{T}\alpha_t h_t(\bm{x}),
  \label{eq:ensemble-additive-model}
\end{equation}
其中$h_t$是第$t$个成员，$\alpha_t$控制其贡献。给定损失$\ell$，理想化的第$t$轮可以写成
\begin{equation}
  (\alpha_t,h_t)
  \in
  \argmin_{\alpha,\,h\in\mathcal{H}}
  \sum_{i=1}^{n}
  \ell\!\left(
    y_i,
    F_{t-1}(\bm{x}_i)+\alpha h(\bm{x}_i)
  \right)
  +\Omega_t(h).
  \label{eq:ensemble-forward-stagewise}
\end{equation}
这里$\mathcal{H}$是成员函数族，$\Omega_t$表示当前方法使用的复杂度约束。前向逐步训练固定已经加入的成员，只优化本轮增量；它避免一次性联合搜索全部函数，却也不等于求得式\eqref{eq:ensemble-additive-model}在所有成员和权重上的全局最优解。

不同Boosting方法的主要差别，可以从“当前不足怎样变成下一轮训练信号”来判断。AdaBoost改变训练样本的权重，使当前错分记录在下一轮具有更大影响；梯度提升把损失对当前预测的负梯度作为伪响应；二阶提升树还利用局部曲率近似叶权重和分裂收益。只有平方损失下的负梯度等于普通响应残差，把所有Boosting概括为“反复拟合残差”会掩盖损失函数的作用。

\subsection{AdaBoost}
\label{sec:ensemble-adaboost}

\term{AdaBoost}（adaptive boosting）从带权训练样本出发，使每一轮弱分类器回应当前加权分布。考虑二分类标签$y_i\in\{-1,1\}$和成员$h_t(\bm{x})\in\{-1,1\}$。初始样本权重为$w_i^{(1)}=1/n$，第$t$轮成员的加权错误率为
\begin{equation}
  \epsilon_t
  =
  \sum_{i=1}^{n}
  w_i^{(t)}
  \mathbb{I}\!\left\{y_i\neq h_t(\bm{x}_i)\right\}.
  \label{eq:ensemble-adaboost-error}
\end{equation}
当$\epsilon_t<1/2$时，成员投票权重取
\begin{equation}
  \alpha_t
  =
  \frac{1}{2}
  \log\frac{1-\epsilon_t}{\epsilon_t},
  \qquad
  w_i^{(t+1)}
  =
  \frac{
    w_i^{(t)}
    \exp\!\left[-\alpha_t y_i h_t(\bm{x}_i)\right]
  }{Z_t},
  \label{eq:ensemble-adaboost-update}
\end{equation}
其中$Z_t$使新权重和为1。错分样本满足$y_i h_t(\bm{x}_i)=-1$，权重乘上$e^{\alpha_t}$；正确样本的权重乘上$e^{-\alpha_t}$。最终分类器按$\operatorname{sign}(\sum_t\alpha_t h_t(\bm{x}))$决策。若成员不能在当前权重下取得小于$1/2$的错误率，标准更新就无法获得正的成员贡献，需要停止、调整成员训练或采用相应变体。

图\ref{fig:ensemble-adaboost-weight-update}展示了式\eqref{eq:ensemble-adaboost-update}形成的反馈回路。当前分类结果先改变样本权重，再由归一化后的权重决定下一轮成员看到的训练分布；错分记录因而获得更大的后续影响。

\begin{figure}[htbp]
  \centering
  % AdaBoost turns classification outcomes into the next round's sample weights.
% Schematic only; circle areas are qualitative and encode no measured values.
% Canvas: 112 mm x 61 mm. Dependencies: project palette and TikZ.
\begin{tikzpicture}[
  x=1mm,y=1mm,
  font=\figurefont\footnotesize,
  text=BookInk,
  stage/.style={
    rectangle,
    model hidden,
    minimum width=24mm,
    minimum height=11mm,
    text width=21mm,
    align=center,
    inner sep=1mm
  },
  data/.style={
    stage,
    model input
  },
  model/.style={
    stage,
    model hidden
  },
  update/.style={
    rectangle,
    model training,
    minimum width=29mm,
    minimum height=11mm,
    text width=26mm,
    align=center,
    inner sep=1mm
  },
  flow/.style={
    model flow
  },
  feedback/.style={
    model feedback
  }
]
  \path[use as bounding box] (0,0) rectangle (112,61);

  \node[anchor=west,font=\figurefont\bfseries\small] at (1,57)
    {第$t$轮：分类结果决定第$t+1$轮的训练分布};

  \node[data] (weighted-data) at (17,35)
    {带权训练样本\\$\{w_i^{(t)}\}$};
  \node[model] (learner) at (49,35)
    {训练弱分类器\\$h_t$};
  \node[update] (outcome) at (91,43)
    {正确分类\\乘以$e^{-\alpha_t}$};
  \node[update] (mistake) at (91,27)
    {错分\\乘以$e^{\alpha_t}$};
  \node[update,minimum width=25mm] (normalize) at (49,10)
    {归一化\\得到$\{w_i^{(t+1)}\}$};

  \draw[flow] (weighted-data)--(learner);
  \draw[flow] (learner.east)--++(4,0)|-(outcome.west);
  \draw[feedback] (learner.east)--++(4,0)|-(mistake.west);
  \draw[flow] (outcome.south)--++(0,-5)-|(normalize.east);
  \draw[feedback] (mistake.south)--++(0,-5)-|(normalize.east);
  \draw[feedback] (normalize.west)--(1,10)|-(weighted-data.west);
  \node[anchor=west,text=NNGradient,align=left] at (33,22)
    {下一轮更关注\\当前错分样本};

  \foreach \x/\r in {6/1.1,12/1.1,18/1.1,24/1.1,30/1.1} {
    \fill[NNInput] (\x,18) circle[radius=\r mm];
  }
  \foreach \x/\r in {6/0.8,12/0.8,18/1.9,24/0.8,30/1.9} {
    \fill[NNGradient] (\x,7) circle[radius=\r mm];
  }
  \node[anchor=west,text=BookMuted] at (2,22)
    {更新前：权重接近};
  \node[anchor=west,text=BookMuted] at (2,2)
    {更新后：错分样本权重变大};
\end{tikzpicture}

  \caption{AdaBoost的样本权重更新。实线表示常规训练与正确分类样本的更新路径，赭色虚线强调错分样本被放大的反馈路径。圆点大小仅示意相对权重变化，不表示具体样本数或实测权重。}
  \label{fig:ensemble-adaboost-weight-update}
\end{figure}

这一更新也可以从指数损失解释。令$F(\bm{x})=\sum_t\alpha_t h_t(\bm{x})$，经验目标为
\[
  \sum_{i=1}^{n}\exp\!\left[-y_iF(\bm{x}_i)\right].
\]
固定$F_{t-1}$后选择$h_t$与$\alpha_t$，可以导出式\eqref{eq:ensemble-adaboost-update}。因此，样本重加权不是与目标函数无关的经验技巧，而是指数损失下前向逐步加法建模的结果\scite{freund1997boosting,friedman2000adaboost}。

指数权重会持续放大难以正确分类的训练记录。若这些记录代表真实而稀少的边界情形，重新分配注意力可能有用；若它们来自标签错误或异常观测，后续成员可能反复追逐不可修正的噪声。训练错误下降还只说明组合在当前样本上的拟合变化，不能替代新样本风险。弱学习条件、训练错误界及投票间隔的泛化分析见第\ref{sec:boosting}节；这里更关心更新如何形成可使用的模型。

多分类AdaBoost需要重新定义成员错误与类别得分，回归变体则要用响应误差决定样本权重和预测聚合。它们保留序贯重加权思想，却不再直接使用上述二分类推导。比较这些变体时，应按第\ref{chap:classification}章或第\ref{chap:regression}章的输出和损失分别评价，不能只沿用二分类准确率。

\subsection{梯度提升机与GBDT}
\label{sec:ensemble-gradient-boosting}

AdaBoost的指数损失给出了特定重加权规则。若任务需要平方损失、对数损失、分位数损失或其他可微目标，可以直接询问：当前预测沿哪个函数方向变化，经验损失下降最快？\term{梯度提升机}（gradient boosting machine, GBM）用基学习器逼近这个负梯度方向，从而把序贯加法扩展到更一般的损失\scite{friedman2001gradientboost}。

在第$t$轮，先对每个训练样本计算\term{伪残差}（pseudo-residual）
\begin{equation}
  r_i^{(t)}
  =
  -
  \left.
  \frac{\partial \ell(y_i,z)}{\partial z}
  \right|_{z=F_{t-1}(\bm{x}_i)}.
  \label{eq:ensemble-pseudo-residual}
\end{equation}
随后训练$h_t$拟合$\{(\bm{x}_i,r_i^{(t)})\}_{i=1}^{n}$，并通过线搜索或叶内优化得到步长$\alpha_t$。加入收缩系数$\eta\in(0,1]$后，
\begin{equation}
  F_t(\bm{x})
  =
  F_{t-1}(\bm{x})
  +
  \eta\alpha_t h_t(\bm{x}).
  \label{eq:ensemble-gradient-boosting-update}
\end{equation}
当$\ell(y,z)=\tfrac12(y-z)^2$时，式\eqref{eq:ensemble-pseudo-residual}给出$y_i-F_{t-1}(\bm{x}_i)$，恰好是普通残差；换成二项对数损失后，伪残差由当前类别概率和标签共同决定，不再等于响应减去任意原始得分。

若$h_t$采用回归树，就得到\term{梯度提升决策树}（gradient boosting decision tree, GBDT）。树把伪残差按叶区域分组，每片叶给出本轮增量；对非平方损失，可以在树结构确定后，再按原损失求各叶的输出值。这里“回归树”描述成员拟合连续修正量的计算角色，整体任务仍可以是分类。

学习率、轮数与单树复杂度共同决定加法函数。较小学习率缩小每轮变化，通常需要更多轮才能达到相近训练目标；较深的树可以在单轮表达更高阶交互，也会扩大每轮容量。行采样和列采样进一步引入随机性，早停则依据独立验证轨迹选择轮数。不存在只调其中一个参数就能固定其他参数作用的普遍规则。

GBM的轮次具有数据依赖，不能像Bagging成员那样全部并行。单轮内部仍可并行统计不同特征或候选分裂。训练损失随轮次下降、优化过程趋于稳定和测试风险达到较小值是三件不同的事；若轮数、损失、树结构和早停规则都根据同一验证集反复调整，还要把这种选择纳入最终评价。

\subsection{XGBoost}
\label{sec:ensemble-xgboost}

经典梯度提升用一阶负梯度给出拟合方向。\term{XGBoost}在正则化加法树目标中进一步使用二阶局部近似，并围绕稀疏输入、近似分裂和内存访问设计训练系统\scite{chen2016xgboost}。这三层需要分开理解：二阶目标决定怎样评价一棵候选树，分裂算法近似搜索树结构，系统机制则决定统计量怎样高效计算。

设第$t$棵树为$f_t$，损失在当前预测$F_{t-1}(\bm{x}_i)$处的一阶导与二阶导分别为$g_i^{(t)}$和$s_i^{(t)}$。去掉与$f_t$无关的常数后，本轮目标近似为
\begin{equation}
  \widetilde{\mathcal{L}}^{(t)}
  =
  \sum_{i=1}^{n}
  \left[
    g_i^{(t)}f_t(\bm{x}_i)
    +\frac{1}{2}s_i^{(t)}f_t^2(\bm{x}_i)
  \right]
  +\gamma L
  +\frac{\lambda}{2}\sum_{j=1}^{L}v_j^2,
  \label{eq:ensemble-xgboost-objective}
\end{equation}
其中$L$是叶数，$v_j$是第$j$片叶的输出，$\gamma$惩罚新增叶，$\lambda$约束叶权重。若$I_j$表示落入第$j$片叶的样本索引，记
\[
  G_j=\sum_{i\in I_j}g_i^{(t)},
  \qquad
  S_j=\sum_{i\in I_j}s_i^{(t)}.
\]
在$S_j+\lambda>0$时，固定树结构的最优叶权重为
\begin{equation}
  v_j^\star=-\frac{G_j}{S_j+\lambda}.
  \label{eq:ensemble-xgboost-leaf-weight}
\end{equation}
把一个父叶分成左右子叶时，相应近似目标下降量为
\begin{equation}
  \operatorname{Gain}
  =
  \frac{1}{2}
  \left[
    \frac{G_{\mathrm{L}}^2}{S_{\mathrm{L}}+\lambda}
    +
    \frac{G_{\mathrm{R}}^2}{S_{\mathrm{R}}+\lambda}
    -
    \frac{(G_{\mathrm{L}}+G_{\mathrm{R}})^2}
         {S_{\mathrm{L}}+S_{\mathrm{R}}+\lambda}
  \right]
  -\gamma.
  \label{eq:ensemble-xgboost-split-gain}
\end{equation}
式\eqref{eq:ensemble-xgboost-split-gain}把当前损失的梯度、曲率与树复杂度放在同一个分裂准则中。它是局部二阶近似下的增益，不表示贪心生长已经找到所有树结构中的全局最优。

图\ref{fig:ensemble-gradient-xgboost}把一阶梯度提升与XGBoost的二阶树评分放在同一计算链中。左图回答新成员拟合什么，右图回答候选树结构怎样利用梯度统计得到叶权重和分裂增益；二阶统计扩展了本轮目标的近似信息，但没有改变逐轮加法更新的外层结构。

\begin{figure*}[htbp]
  \centering
  % Gradient boosting's first-order signal and XGBoost's second-order tree score.
% Schematic only; no empirical quantities are encoded.
% Canvas: 164 mm x 64 mm. Dependencies: project palette and TikZ.
\begin{tikzpicture}[
  x=1mm,y=1mm,
  font=\footnotesize,
  text=BookInk,
  state/.style={
    rectangle,
    draw=FigureInput,
    fill=FigureInput!8,
    line width=0.7pt,
    minimum width=26mm,
    minimum height=10mm,
    text width=22mm,
    align=center,
    inner sep=1mm
  },
  model/.style={
    rectangle,
    draw=FigureModel,
    fill=FigureModel!8,
    line width=0.8pt,
    minimum width=28mm,
    minimum height=10mm,
    text width=24mm,
    align=center,
    inner sep=1mm
  },
  signal/.style={
    rectangle,
    draw=FigureLoss,
    fill=FigureLoss!9,
    line width=0.8pt,
    minimum width=30mm,
    minimum height=11mm,
    text width=26mm,
    align=center,
    inner sep=1mm
  },
  result/.style={
    rectangle,
    draw=FigureOutput,
    fill=FigureOutput!6,
    line width=0.8pt,
    minimum width=29mm,
    minimum height=11mm,
    text width=25mm,
    align=center,
    inner sep=1mm
  },
  flow/.style={
    -{Latex[length=2.2mm,width=1.5mm]},
    draw=BookMuted,
    line width=0.8pt,
    rounded corners=1.2mm
  }
]
  \path[use as bounding box] (0,0) rectangle (164,64);

  \node[anchor=west,font=\heiti\small] at (1,60)
    {(a) 梯度提升：把损失方向变成新成员};
  \node[state] (current) at (18,42)
    {当前预测\\$F_{t-1}(\bm{x}_i)$};
  \node[signal] (gradient) at (52,42)
    {计算负梯度\\$r_i^{(t)}=-\partial_z\ell$};
  \node[model] (fit) at (52,19)
    {用$h_t$拟合\\伪残差};
  \node[result] (updated) at (18,19)
    {加法更新\\$F_t=F_{t-1}+\eta\alpha_t h_t$};
  \draw[flow] (current)--(gradient);
  \draw[flow] (gradient)--(fit);
  \draw[flow] (fit)--(updated);
  \draw[flow] (updated)--(current);

  \draw[BookRule,line width=0.6pt] (78,3)--(78,61);

  \node[anchor=west,font=\heiti\small] at (83,60)
    {(b) XGBoost：二阶统计进入树分裂};
  \node[signal] (derivatives) at (101,43)
    {样本统计\\$(g_i^{(t)},s_i^{(t)})$};
  \node[model] (partition) at (139,43)
    {候选分裂\\$I\to(I_{\mathrm L},I_{\mathrm R})$};
  \node[signal] (sums) at (101,21)
    {按候选叶聚合\\$(G_j,S_j)$};
  \node[result] (score) at (139,21)
    {叶权重$v_j^\star$\\与分裂增益};
  \node[signal,minimum width=29mm,minimum height=7mm,text width=25mm]
    (penalty) at (139,6)
    {复杂度惩罚$(\gamma,\lambda)$};
  \draw[flow] (derivatives)--(sums);
  \draw[flow] (partition.south)--++(0,-6)-|(sums.east);
  \draw[flow] (sums)--(score);
  \draw[flow] (penalty)--(score);
\end{tikzpicture}

  \caption{梯度提升与XGBoost中的本轮修正。左图中，损失对当前预测的负梯度形成伪残差，新成员拟合该信号后更新组合；右图中，XGBoost在候选叶内汇总一阶导与二阶导，并结合复杂度惩罚计算叶权重和分裂增益。该图描述局部优化机制，不表示贪心树结构是全局最优。}
  \label{fig:ensemble-gradient-xgboost}
\end{figure*}

精确分裂搜索遍历排序后的候选特征值；近似搜索则先用带权分位数草图构造候选切分，再汇总梯度统计。稀疏感知搜索为缺失值学习默认方向，列采样减少每轮需要检查的特征。块式存储、缓存感知访问、并行统计和外存计算改变的是实现代价，不改变式\eqref{eq:ensemble-xgboost-objective}所定义的学习对象。工程改进能否转化为当前硬件和预算下的速度收益，需要实测。

\subsection{LightGBM}
\label{sec:ensemble-lightgbm}

当样本和特征很多时，逐值搜索分裂会反复扫描大量候选。\term{LightGBM}以直方图训练为基础，把连续特征映射到有限个桶，并在桶级别累计梯度和曲率。构造直方图仍要访问节点样本，但候选分裂只需扫描桶；已知父节点和一个子节点的直方图时，还可以用差分得到另一个子节点的统计量。分桶减少存储和候选数，也引入离散化误差。

树的生长顺序进一步改变计算分配。按层生长同时扩展某一深度上的叶，按叶生长则每次选择当前增益最大的叶继续分裂。在相同叶数下，后者把容量集中到训练目标下降较大的局部区域，可能形成深度不均衡的树。叶数、叶内最小样本量和最大深度因此需要联合约束；训练损失下降更快不等于测试风险必然更小。

LightGBM论文进一步提出两种针对大规模数据的近似。\term{单边梯度采样}（gradient-based one-side sampling, GOSS）保留绝对梯度较大的样本，从绝对梯度较小的样本中随机抽取一部分，并在分裂统计中重新加权后者；它近似的是用于选择分裂的数据统计。\term{互斥特征捆绑}（exclusive feature bundling, EFB）把很少同时非零的稀疏特征编码到较少的组合特征中，近似减少有效特征数。GOSS、EFB和直方图分别作用于样本、特征与候选阈值，不能合并成一个笼统的“降采样”步骤\scite{ke2017lightgbm}。

图\ref{fig:ensemble-lightgbm-approximations}按近似对象并列这三种机制。直方图压缩候选阈值，GOSS压缩参与分裂统计的样本，EFB压缩稀疏特征表示；读者据此可以分别追踪每种近似改变了哪一项计算。

\begin{figure*}[htbp]
  \centering
  % LightGBM reduces split-search work along threshold, sample, and feature axes.
% Schematic only; marks and sparsity patterns encode no measured values.
% Canvas: 164 mm x 64 mm. Dependencies: project palette and TikZ.
\begin{tikzpicture}[
  x=1mm,y=1mm,
  font=\footnotesize,
  text=BookInk,
  flow/.style={
    -{Latex[length=2.2mm,width=1.5mm]},
    draw=BookMuted,
    line width=0.8pt
  },
  keep/.style={
    circle,
    draw=FigureLoss,
    fill=FigureLoss!12,
    line width=0.8pt,
    minimum size=5mm,
    inner sep=0pt
  },
  sample/.style={
    circle,
    draw=FigureInput,
    fill=FigureInput!8,
    line width=0.7pt,
    minimum size=4mm,
    inner sep=0pt
  }
]
  \path[use as bounding box] (0,0) rectangle (164,64);

  \node[font=\heiti\small] at (27,59) {(a) 直方图：减少候选阈值};
  \draw[BookMuted,line width=0.7pt] (5,35)--(49,35);
  \foreach \x in {8,12,15,18,22,24,28,31,35,39,42,46} {
    \fill[FigureInput] (\x,35) circle[radius=1.1mm];
  }
  \foreach \x in {5,16,27,38,49} {
    \draw[FigureModel,line width=0.8pt] (\x,28)--(\x,43);
  }
  \node[text=BookMuted] at (27,24) {连续取值映射到有限个桶};
  \node[draw=FigureOutput,fill=FigureOutput!6,line width=0.8pt,
    minimum width=38mm,minimum height=9mm,align=center] (bins) at (27,10)
    {只扫描桶边界\\并累计梯度统计};
  \draw[flow] (27,22)--(bins);

  \draw[BookRule,line width=0.6pt] (54.5,3)--(54.5,61);

  \node[font=\heiti\small] at (82,59) {(b) GOSS：减少参与统计的样本};
  \node[anchor=east,text=BookMuted] at (72,47) {较大的$|g_i|$};
  \foreach \x in {78,86,94,102} {
    \node[keep] at (\x,47) {};
  }
  \node[anchor=east,text=BookMuted] at (72,34) {较小的$|g_i|$};
  \foreach \x in {75,81,87,93,99,105} {
    \node[sample] at (\x,34) {};
  }
  \node[draw=FigureOutput,fill=FigureOutput!6,line width=0.8pt,
    minimum width=42mm,minimum height=15mm,text width=38mm,align=center]
    (goss) at (82,12)
    {保留$|g_i|$较大的样本\\随机抽取并重加权部分$|g_i|$较小的样本};
  \draw[flow,draw=FigureLoss] (105,47)--(108,47)|-(goss.east);
  \draw[flow] (87,31.5)--(87,27)--(57,27)|-(goss.west);
  \draw[BookMuted,dashed,line width=0.7pt] (81,31)--(81,37);
  \draw[BookMuted,dashed,line width=0.7pt] (93,31)--(93,37);
  \draw[BookMuted,dashed,line width=0.7pt] (105,31)--(105,37);

  \draw[BookRule,line width=0.6pt] (109.5,3)--(109.5,61);

  \node[font=\heiti\small] at (137,59) {(c) EFB：减少有效特征};
  \node[anchor=east,text=BookMuted] at (121,45) {$x_a$};
  \node[anchor=east,text=BookMuted] at (121,35) {$x_b$};
  \foreach \x in {125,131,137,143,149,155} {
    \draw[BookRule,line width=0.6pt] (\x-2.2,42.5) rectangle (\x+2.2,47.5);
    \draw[BookRule,line width=0.6pt] (\x-2.2,32.5) rectangle (\x+2.2,37.5);
  }
  \fill[FigureInput] (122.8,42.5) rectangle (127.2,47.5);
  \fill[FigureInput] (140.8,42.5) rectangle (145.2,47.5);
  \fill[FigureInput] (128.8,32.5) rectangle (133.2,37.5);
  \fill[FigureInput] (146.8,32.5) rectangle (151.2,37.5);
  \node[text=BookMuted,align=center] at (139,25)
    {非零位置很少冲突};
  \node[draw=FigureOutput,fill=FigureOutput!6,line width=0.8pt,
    minimum width=42mm,minimum height=10mm,text width=38mm,align=center]
    (bundle) at (137,10)
    {编码为一个捆绑特征};
  \draw[flow] (139,21)--(bundle);
\end{tikzpicture}

  \caption{LightGBM中三种作用位置不同的计算近似。直方图把连续取值映射到有限桶，GOSS保留$|g_i|$较大的样本，并抽取、重加权部分$|g_i|$较小的样本，EFB把很少同时非零的特征编码为较少的捆绑特征。图中点数、桶数与非零位置均为机制示意。}
  \label{fig:ensemble-lightgbm-approximations}
\end{figure*}

这些机制的效果依赖梯度分布、特征稀疏性、桶数、冲突处理和硬件实现。论文中的速度与精度比较支持相应实验条件下的结论，不构成LightGBM相对其他实现的无条件排序。公平比较应固定损失、搜索预算和早停规则，并记录离散化与采样设置。

\subsection{CatBoost}
\label{sec:ensemble-catboost}

高基数类别特征若全部展开为独热变量，会增加特征规模；把类别直接映射为整数，又会引入原本不存在的次序。另一种做法是用同类别样本的响应均值编码，但若样本自己的$y_i$参与其编码，模型就能从输入中间接看到待预测标签。这种\term{目标泄漏}（target leakage）会使训练特征与部署特征的构造条件不同。

\term{CatBoost}用随机排列建立\term{有序目标统计}（ordered target statistic）。设$c_i$为样本$i$的某个类别值，用$\pi(j)$表示样本$j$在随机排列中的位置；对当前样本只使用排列中位于它之前且类别相同的记录：
\begin{equation}
  z_i
  =
  \frac{
    \sum_{j:\,\pi(j)<\pi(i)}
    \mathbb{I}\{c_j=c_i\}y_j
    +ap
  }{
    \sum_{j:\,\pi(j)<\pi(i)}
    \mathbb{I}\{c_j=c_i\}
    +a
  }.
  \label{eq:ensemble-catboost-ordered-statistic}
\end{equation}
其中$p$是先验统计量，$a>0$控制平滑。样本自身的标签不会进入$z_i$；部署时则可使用训练集中的类别统计。低频类别仍会产生较大估计波动，平滑、多个排列和类别组合需要共同考虑。

图\ref{fig:ensemble-catboost-ordered-statistics}用一个排列位置说明这种隔离。计算当前类别$A$样本的编码时，只允许排列前缀中同为类别$A$的标签进入统计量；其他类别、当前标签和排列后缀都不参与这次编码。

\begin{figure}[htbp]
  \centering
  % CatBoost ordered target statistic for one sample in a random permutation.
% Schematic only; category symbols and labels are placeholders.
% Canvas: 112 mm x 58 mm. Dependencies: project palette and TikZ.
\begin{tikzpicture}[
  x=1mm,y=1mm,
  font=\footnotesize,
  text=BookInk,
  record/.style={
    rectangle,
    draw=FigureInput,
    fill=FigureInput!8,
    line width=0.7pt,
    minimum width=17mm,
    minimum height=14mm,
    text width=14mm,
    align=center,
    inner sep=1mm
  },
  current/.style={
    record,
    draw=FigureLoss,
    fill=FigureLoss!10,
    line width=1pt
  },
  future/.style={
    record,
    draw=BookRule,
    fill=BookPaper,
    text=BookMuted
  },
  result/.style={
    rectangle,
    draw=FigureOutput,
    fill=FigureOutput!6,
    line width=0.8pt,
    minimum width=42mm,
    minimum height=12mm,
    text width=38mm,
    align=center,
    inner sep=1mm
  },
  flow/.style={
    -{Latex[length=2.2mm,width=1.5mm]},
    draw=BookMuted,
    line width=0.8pt,
    rounded corners=1.2mm
  },
  use/.style={
    -{Latex[length=2.2mm,width=1.5mm]},
    draw=BookMuted,
    line width=0.9pt,
    rounded corners=1.2mm
  }
]
  \path[use as bounding box] (0,0) rectangle (112,58);

  \node[anchor=west,font=\heiti\small] at (1,54)
    {按排列位置$\pi(j)$从小到大的训练样本};
  \draw[flow] (4,47)--(108,47);

  \node[record] (p1) at (14,32)
    {$j_1$\\$\pi(j_1)<\pi(i)$\\$c_{j_1}=A$\\$y_{j_1}$};
  \node[record] (p2) at (36,32)
    {$j_2$\\$\pi(j_2)<\pi(i)$\\$c_{j_2}=B$\\$y_{j_2}$};
  \node[record] (p3) at (58,32)
    {$j_3$\\$\pi(j_3)<\pi(i)$\\$c_{j_3}=A$\\$y_{j_3}$};
  \node[current] (pi) at (80,32)
    {样本$i$\\$\pi(i)$\\$c_i=A$\\不使用$y_i$};
  \node[future] (pf) at (102,32)
    {$j_f$\\$\pi(j_f)>\pi(i)$\\即使$c_{j_f}=A$\\也不使用};
  \foreach \name in {p1,p2,p3,pi,pf} {
    \draw[BookMuted,line width=0.7pt] (\name.north)--++(0,6);
  }

  \node[result] (statistic) at (47,10)
    {只汇总$\pi(j)<\pi(i)$且$c_j=c_i$的标签$y_j$\\
     得到样本$i$的编码$z_i$};
  \draw[use] (p1.south)--++(0,-5)-|(statistic.north);
  \draw[use] (p3.south)--++(0,-5)-|(statistic.north);
  \draw[BookMuted,dashed,line width=0.8pt] (p2.south)--++(0,-5);
  \draw[BookMuted,dashed,line width=0.8pt] (pi.south)--++(0,-5);
  \draw[BookMuted,dashed,line width=0.8pt] (pf.south)--++(0,-5);
\end{tikzpicture}

  \caption{CatBoost有序目标统计的前缀约束。这里$\pi(j)$表示样本$j$在随机排列中的位置；计算当前样本$i$的编码时，指向统计量的实线只连接满足$\pi(j)<\pi(i)$且$c_j=c_i$的标签。灰色虚线表示其他类别、当前标签和后续记录在本次统计中被隔离。样本索引与类别符号仅用于说明机制。}
  \label{fig:ensemble-catboost-ordered-statistics}
\end{figure}

有序提升把相同的前缀思想用于梯度估计。直观地说，计算某个训练样本的当前预测与伪残差时，只使用排列中位于该样本之前的记录形成相应模型，减少“同一样本既参与拟合当前模型，又用于评价当前模型”的预测偏移。实际算法通过多个排列和前缀模型实现这一约束，训练状态与计算量也比普通单序列提升更复杂\scite{prokhorenkova2018catboost}。

CatBoost常使用\term{对称树}（symmetric tree），即同一深度的所有节点采用相同分裂条件。深度为$d_{\mathrm{tree}}$的完整对称树至多有$2^{d_{\mathrm{tree}}}$片叶，路径可以编码为固定长度的二进制选择。这种结构便于批量预测，也限制了单棵树可表达的分区形状。类别编码、有序提升和对称树分别处理输入构造、序贯估计偏移和成员结构，不能把某一项的作用归给整个系统。

\subsection{模型对比与总结}
\label{sec:ensemble-boosting-comparison}

表\ref{tab:ensemble-boosting-comparison}先固定共同的序贯加法结构，再比较各方法修改哪一层。AdaBoost由加权分类错误生成样本权重；GBM由损失梯度生成伪残差；XGBoost对正则化树目标作二阶近似；LightGBM进一步压缩分裂统计的计算；CatBoost重点处理类别统计与序贯估计中的目标泄漏。后四者之间存在实现重叠，表中只列代表性机制，不把软件当前支持的全部选项写成永恒差异。

\begin{table*}[htbp]
  \centering
  \small
  \begin{tabularx}{\linewidth}{@{}>{\raggedright\arraybackslash}p{19mm}*{6}{>{\raggedright\arraybackslash}X}@{}}
    \toprule
    方法 & 本轮修正信号 & 典型成员 & 目标近似 & 树生长或分裂 & 类别特征重点 & 主要约束 \\
    \midrule
    AdaBoost & 加权错分 & 树桩或浅分类树 & 指数损失下的解析更新 & 由成员算法决定 & 需另行编码 & 对持续错分记录敏感 \\
    GBM／GBDT & 负梯度伪残差 & 回归树 & 一阶函数方向与叶内优化 & 通常贪心树生长 & 需随实现处理 & 学习率、轮数与树复杂度耦合 \\
    XGBoost & 一阶导与二阶导 & 正则化回归树 & 二阶局部目标 & 精确或分位数候选；稀疏默认方向 & 依实现与版本而定 & 轮次串行，近似与系统配置影响成本 \\
    LightGBM & 桶级梯度统计 & 直方图回归树 & 通常使用二阶统计 & 按叶生长；直方图候选 & 支持类别分组，机制依实现 & 分桶、GOSS与EFB引入近似条件 \\
    CatBoost & 有序统计与有序梯度 & 对称树 & 梯度或二阶统计 & 同层共享分裂条件 & 有序目标统计与组合 & 多排列和前缀状态增加训练成本 \\
    \bottomrule
  \end{tabularx}
  \caption{Boosting代表方法的机制比较。GBM是一般框架，GBDT指使用回归树成员的实例；XGBoost、LightGBM和CatBoost是树提升的具体模型与系统，不能与Boosting这一上位类别并列。}
  \label{tab:ensemble-boosting-comparison}
\end{table*}

随机梯度提升在每轮使用训练子样本，DART随机暂时移除部分已有树后计算修正，LogitBoost则从加法Logistic回归推导更新\scite{friedman2002stochastic,rashmi2015dart,friedman2000adaboost}。理解这些扩展时，可以询问它改变了训练样本、已有成员、损失近似还是树结构。Boosting能够沿训练目标定向增加容量，却仍采用预先规定的加法组合形式；若已有成员来自不同模型族，并希望从它们的样本外表现中学习组合方式，就进入Stacking的范围。

\section{基于Stacking的集成}
\label{sec:ensemble-stacking}

Bagging预先规定平均或投票，Boosting在训练中形成加法权重；二者的成员生成机制与组合方式紧密绑定。若已经有线性模型、近邻模型、随机森林和梯度提升树等异质候选，还可以把它们的预测视为新特征，再训练一个模型决定怎样组合。\term{Stacking}（stacked generalization，堆叠）采用这种两层结构：基学习器产生一级预测，\term{元学习器}（meta-learner）把这些预测映射为最终输出\scite{wolpert1992stacked,breiman1996stacked}。

\subsection{Stacking}
\label{sec:ensemble-stacking-model}

\subsubsection{为什么不能直接使用训练内预测}

设第$m$个成员模型为$h_m$。若直接在完整训练集上拟合$h_m$，再把$h_m(\bm{x}_i)$作为样本$i$的元特征，元学习器看到的是成员对自身训练数据的预测。一个容量很大的成员可能近乎记住训练标签，元学习器便会高估它在新样本上的可靠性。问题不在于元学习器“多看了一列数据”，而在于这列特征的生成条件与部署时不同。

Stacking通常用交叉验证生成\term{折外预测}（out-of-fold prediction, OOF prediction）。把训练样本划分为$K$个互不相交的折$V_1,\ldots,V_K$。对每个成员类型$m$和每个折$k$，只用$D_n\setminus V_k$训练临时模型$h_m^{(-k)}$，再预测$V_k$中的样本。若样本$i$属于$V_{k(i)}$，元特征定义为
\begin{equation}
  z_{i,m}
  =
  h_m^{(-k(i))}(\bm{x}_i).
  \label{eq:ensemble-oof-feature}
\end{equation}
把所有成员的折外预测拼成$\bm{z}_i=(z_{i,1},\ldots,z_{i,M})^{\top}$，再由元学习算法$\mathcal{G}$训练
\[
  g=\mathcal{G}\!\left(\{(\bm{z}_i,y_i)\}_{i=1}^{n}\right).
\]
这样，每个$z_{i,m}$都来自未使用样本$i$标签拟合的临时成员。图\ref{fig:ensemble-stacking-oof}展示了折外矩阵的形成过程。

\begin{figure*}[htbp]
  \centering
  % Three-fold out-of-fold feature construction for Stacking.
% Schematic only; no empirical quantities are encoded.
% Canvas: 164 mm x 68 mm. Dependencies: project palette and TikZ.
\begin{tikzpicture}[
  x=1mm,y=1mm,
  font=\figurefont\footnotesize,
  text=BookInk,
  fold/.style={
    rectangle,
    model input,
    minimum width=16mm,
    minimum height=9mm,
    align=center,
    inner sep=1mm
  },
  train/.style={
    rectangle,
    model context,
    minimum width=27mm,
    minimum height=9mm,
    align=center,
    inner sep=1mm
  },
  model/.style={
    rectangle,
    model hidden,
    minimum width=25mm,
    minimum height=9mm,
    align=center,
    inner sep=1mm
  },
  predict/.style={
    rectangle,
    model training,
    minimum width=25mm,
    minimum height=9mm,
    align=center,
    inner sep=1mm
  },
  matrix/.style={
    rectangle,
    model output,
    minimum width=18mm,
    minimum height=43mm,
    text width=14mm,
    align=center,
    inner sep=1mm
  },
  meta/.style={
    rectangle,
    model output,
    minimum width=24mm,
    minimum height=9mm,
    align=center,
    inner sep=1mm
  },
  flow/.style={
    model flow
  }
]
  \path[use as bounding box] (0,0) rectangle (164,68);

  \node[anchor=west,font=\figurefont\bfseries\small] at (1,64) {三折轮换};

  \node[fold] (v1) at (12,52) {留出$V_1$};
  \node[fold] (v2) at (12,35) {留出$V_2$};
  \node[fold] (v3) at (12,18) {留出$V_3$};

  \node[train] (t1) at (42,52) {训练$D_n\setminus V_1$};
  \node[train] (t2) at (42,35) {训练$D_n\setminus V_2$};
  \node[train] (t3) at (42,18) {训练$D_n\setminus V_3$};

  \node[model] (m1) at (75,52) {$h_1^{(-1)},\ldots,h_M^{(-1)}$};
  \node[model] (m2) at (75,35) {$h_1^{(-2)},\ldots,h_M^{(-2)}$};
  \node[model] (m3) at (75,18) {$h_1^{(-3)},\ldots,h_M^{(-3)}$};

  \node[predict] (p1) at (108,52) {预测$V_1$};
  \node[predict] (p2) at (108,35) {预测$V_2$};
  \node[predict] (p3) at (108,18) {预测$V_3$};

  \node[matrix] (z) at (145,35)
    {按原样本顺序拼接\\[1mm]
     折外矩阵\\$\bm{Z}\in\mathbb{R}^{n\times M}$};
  \node[meta] (g) at (145,5.5) {元学习器$g$};

  \draw[flow] (v1) -- (t1);
  \draw[flow] (t1) -- (m1);
  \draw[flow] (v2) -- (t2);
  \draw[flow] (t2) -- (m2);
  \draw[flow] (v3) -- (t3);
  \draw[flow] (t3) -- (m3);
  \draw[flow] (m1) -- (p1);
  \draw[flow] (m2) -- (p2);
  \draw[flow] (m3) -- (p3);
  \draw[flow] (p1.east) -- ([yshift=17mm]z.west);
  \draw[flow] (p2.east) -- (z.west);
  \draw[flow] (p3.east) -- ([yshift=-17mm]z.west);
  \draw[flow] (z.south) -- (g.north);
\end{tikzpicture}

  \caption{三折Stacking的折外特征构造。每一折轮流作为留出折，临时成员只在其余两折训练，并为留出折产生预测；拼接三次预测得到完整折外矩阵$\bm{Z}$，元学习器以$\bm{Z}$为输入并使用对应标签拟合。部署阶段的基成员还需按预先规定的方式重训或保留。}
  \label{fig:ensemble-stacking-oof}
\end{figure*}

分类成员若输出$C$维概率向量，一个成员通常贡献$C$列元特征；二分类时也可以只保留一个类别概率。硬标签丢失了置信程度，未经校准的得分又可能量纲不同。采用哪种输出必须在训练和部署阶段保持一致，元学习器的正则化也要适应该尺度与相关性。

\subsubsection{部署模型与完整数据重训}

折外矩阵用于训练元学习器，但部署时还需要能预测新输入的基成员。一种做法是在完整训练集上重新拟合每类基学习器$\widehat{h}_m$，预测
\[
  \widehat{y}
  =
  g\!\left(
    \widehat{h}_1(\bm{x}),\ldots,\widehat{h}_M(\bm{x})
  \right).
\]
这种做法充分使用训练数据，却使元学习器训练时看到的折内模型预测与部署时的全量模型预测存在有限样本差异。另一种做法是保留每折模型，对新输入先在折内模型间平均，再交给元学习器；其输入条件更接近折外构造，但需要保存和运行更多模型。

两种部署方案都可以成立，关键是事先固定并完整评价实际链路。基学习器的特征处理、类别编码和超参数搜索必须在相应训练折内完成。若先用全部数据选择特征或编码，再只把模型拟合放进交叉验证，信息仍可能从留出折进入元特征。

\subsection{Blending}
\label{sec:ensemble-blending}

\term{Blending}用一次固定留出替代$K$折交叉验证。将可用于开发的数据分为基训练集$D_{\mathrm{base}}$和组合训练集$D_{\mathrm{blend}}$；基成员只在$D_{\mathrm{base}}$上拟合，对$D_{\mathrm{blend}}$产生预测，元学习器再以这些预测和对应标签训练。每个基成员只需完成一次用于元特征的拟合，流程比完整Stacking短。

代价是元学习器只能使用$D_{\mathrm{blend}}$中的样本，基成员也暂时失去这部分训练数据。划分较小时，元学习器估计不稳定；划分较大时，基成员训练信息减少。多次随机划分并平均可以减小单次划分的偶然性，但它重新引入多次训练，已经不再具有单次Blending的全部成本优势。

部署时若把$D_{\mathrm{blend}}$加入基成员重新训练，基预测分布可能相对元学习器训练时发生变化；若不加入，则没有充分利用已有标签。这个取舍不能由“Blending更简单”自动解决，需要在流程中明确重训方式，并用冻结后的验证或测试数据评价完整组合。

\subsection{受约束线性Stacking}
\label{sec:ensemble-constrained-linear-stacking}

元学习器不一定要很复杂。对回归输出，可以在折外矩阵$\bm{Z}\in\mathbb{R}^{n\times M}$上学习线性组合
\begin{equation}
  (\widehat{b},\widehat{\bm{w}})
  \in
  \argmin_{b,\bm{w}}
  \frac{1}{n}
  \sum_{i=1}^{n}
  \ell\!\left(
    y_i,
    b+\bm{w}^{\top}\bm{z}_i
  \right)
  +\lambda\lVert\bm{w}\rVert_2^2.
  \label{eq:ensemble-linear-stacking}
\end{equation}
也可以施加$w_m\geq0$或$\sum_mw_m=1$，把组合限制为成员预测的凸组合。非负和归一化约束减少高度相关成员之间的大幅正负抵消，使结果不易外推到成员预测范围之外；但若真实有效的校正需要负权重或截距，这些约束也会增加偏差。

权重大小同时受输出尺度、成员相关性和损失函数影响。两个成员预测近乎相同，线性模型可以在它们之间任意分配相近作用；一个成员输出的数值范围更大，也会改变未标准化正则项下的权重。因此，$\widehat{w}_m$不能直接解释为成员的独立重要性。对类别概率做凸组合时，结果仍位于概率单纯形内，却仍需检查样本外校准。

\subsection{Super Learner}
\label{sec:ensemble-super-learner}

\term{Super Learner}把候选学习器库、交叉验证风险与组合权重放进一个明确程序。给定预先规定的候选库和损失，先生成每个候选的折外预测，再在指定组合类中选择交叉验证风险较小的组合；组合类可以是最佳单成员、非负线性组合或其他受控元模型。

在独立同分布、损失有界或满足相应矩条件、候选库与交叉验证方案受控等条件下，Super Learner的理论结果把其风险与候选组合类中的oracle风险联系起来：样本量增加时，交叉验证选择的组合能够渐近接近该库中最优候选或最优组合的风险\scite{vanderlaan2007superlearner}。这里的oracle只在预先给定的库和损失内比较，不表示找到了所有可测预测器中的最优模型。

三个条件不能混为同一句“组合必然优于成员”。候选库中需要存在有用且互补的预测；折外风险需要在当前数据制度下稳定估计；元学习器还要有足够样本控制自身复杂度。渐近oracle结论也不推出任意有限样本上都优于最佳成员。若数据存在时间、群组或空间依赖，交叉验证划分还必须与部署时的新样本定义一致。

\subsection{模型对比与总结}
\label{sec:ensemble-stacking-comparison}

表\ref{tab:ensemble-stacking-comparison}比较Stacking家族的元特征来源与组合约束。标准Stacking和Super Learner都依赖折外预测；Blending以一次留出降低训练次数；受约束线性Stacking说明“学习组合”不等于必须使用高容量元模型。模型名称之外，还应记录预处理、基模型重训和元输出类型，否则无法复现实际预测链路。

\begin{table*}[htbp]
  \centering
  \small
  \begin{tabularx}{\linewidth}{@{}>{\raggedright\arraybackslash}p{25mm}*{5}{>{\raggedright\arraybackslash}X}@{}}
    \toprule
    方法 & 元特征来源 & 数据利用 & 组合规则 & 主要成本 & 主要风险 \\
    \midrule
    Stacking & $K$折折外预测 & 每个样本都产生元特征 & 可线性或非线性 & 每类成员需多次折内拟合 & 折外流程不完整造成泄漏 \\
    Blending & 固定留出集预测 & 元学习器只用留出集 & 可线性或非线性 & 基成员用于元特征时只拟合一次 & 划分敏感与重训分布变化 \\
    受约束线性Stacking & 通常为折外预测 & 与所用划分一致 & 非负、和为一或正则化线性组合 & 组合层较轻量 & 约束可能排除有效校正 \\
    Super Learner & 交叉验证预测与风险 & 由交叉验证方案决定 & 在预定库和组合类中选择 & 候选库越大，训练与选择成本越高 & oracle范围和有限样本被夸大 \\
    \bottomrule
  \end{tabularx}
  \caption{Stacking及其代表性变体。表中“数据利用”描述元学习器训练阶段；部署前是否用完整开发数据重训基成员，还需作为流程的一部分单独规定。}
  \label{tab:ensemble-stacking-comparison}
\end{table*}

多层Stacking把一层折外输出继续交给下一层，每增加一层都要重新保证当前样本没有通过上游拟合泄漏到本层特征。输入相关的门控组合则让权重随$\bm{x}$变化，它与第\ref{sec:ensemble-dynamic-selection}节的动态成员选择相邻，也与神经网络中的混合专家相接。它们增加了表达能力，同时扩大了数据划分、模型存储和调试成本；在两层组合尚未提供稳定新增价值时，没有理由仅凭层数增加继续扩展。

\section{成员选择与集成简化}
\label{sec:ensemble-selection-compression}

Bagging、Boosting和Stacking先回答怎样形成组合。训练完成后，成员数量仍可能超过存储或延迟预算，也可能包含只在少数输入区域有用的模型。简化集成有三种不同目标：为所有输入固定一个较小成员子集，针对当前输入临时选择成员，或训练一个新的紧凑模型替代原集成。前两种方法保留部分原成员，第三种方法改变了部署模型。

图\ref{fig:ensemble-selection-compression}沿部署对象区分这三种目标。静态选择把成员池缩成一个固定子集，动态选择保留成员池并让路由依赖当前输入，紧凑化则把完整集成的行为转移到新模型。后续评价资源收益时，需要测量图中实际部署的整条路径。

\begin{figure*}[htbp]
  \centering
  % Static selection, dynamic selection, and compression change deployment differently.
% Schematic only; no empirical quantities are encoded.
% Canvas: 164 mm x 64 mm. Dependencies: project palette and TikZ.
\begin{tikzpicture}[
  x=1mm,y=1mm,
  font=\footnotesize,
  text=BookInk,
  pool/.style={
    rectangle,
    draw=FigureModel,
    fill=FigureModel!8,
    line width=0.7pt,
    minimum width=42mm,
    minimum height=10mm,
    align=center,
    inner sep=1mm
  },
  process/.style={
    rectangle,
    draw=FigureLoss,
    fill=FigureLoss!9,
    line width=0.8pt,
    minimum width=42mm,
    minimum height=11mm,
    text width=38mm,
    align=center,
    inner sep=1mm
  },
  deployed/.style={
    rectangle,
    draw=FigureOutput,
    fill=FigureOutput!6,
    line width=0.8pt,
    minimum width=42mm,
    minimum height=12mm,
    text width=38mm,
    align=center,
    inner sep=1mm
  },
  flow/.style={
    -{Latex[length=2.2mm,width=1.5mm]},
    draw=BookMuted,
    line width=0.8pt
  }
]
  \path[use as bounding box] (0,0) rectangle (164,64);

  \node[font=\heiti\small] at (27,60) {(a) 静态成员选择};
  \node[pool] (static-pool) at (27,47)
    {已训练成员池\\$h_1,\ldots,h_M$};
  \node[process] (static-select) at (27,29)
    {依据统一验证目标\\固定子集$S$};
  \node[deployed] (static-deploy) at (27,10)
    {所有输入运行\\同一个成员子集};
  \draw[flow] (static-pool)--(static-select);
  \draw[flow] (static-select)--(static-deploy);

  \node[font=\heiti\small] at (82,60) {(b) 动态成员选择};
  \node[pool] (dynamic-pool) at (82,47)
    {已训练成员池\\$h_1,\ldots,h_M$};
  \node[process] (dynamic-select) at (82,29)
    {根据查询输入$\bm{x}$\\估计局部能力};
  \node[deployed] (dynamic-deploy) at (82,10)
    {不同输入可运行\\不同成员或权重};
  \draw[flow] (dynamic-pool)--(dynamic-select);
  \draw[flow] (dynamic-select)--(dynamic-deploy);

  \node[font=\heiti\small] at (137,60) {(c) 紧凑化};
  \node[pool] (teacher) at (137,47)
    {完整集成教师\\$H$};
  \node[process] (distill) at (137,29)
    {用标签与教师输出\\训练学生};
  \node[deployed] (student) at (137,10)
    {部署新的紧凑模型\\$s_{\bm{\theta}}$};
  \draw[flow] (teacher)--(distill);
  \draw[flow] (distill)--(student);

  \draw[BookRule,line width=0.6pt] (54.5,2)--(54.5,61);
  \draw[BookRule,line width=0.6pt] (109.5,2)--(109.5,61);
\end{tikzpicture}

  \caption{静态成员选择、动态成员选择与紧凑化改变的部署对象。静态选择为所有输入保留同一成员子集；动态选择根据输入改变运行成员或组合权重；紧凑化训练并部署新的学生模型。三种路径可能减少不同类型的成本，不能只用保留成员数比较。}
  \label{fig:ensemble-selection-compression}
\end{figure*}

\subsection{静态成员选择}
\label{sec:ensemble-static-selection}

\term{静态成员选择}（static ensemble selection）给定已经训练好的候选池$\{h_1,\ldots,h_M\}$，选择一个对所有未来输入都相同的子集$S\subseteq\{1,\ldots,M\}$。若$\mathcal{L}_{\mathrm{val}}(S)$表示该子集按指定规则组合后的验证损失，可以写成
\begin{equation}
  \widehat{S}
  \in
  \argmin_{S\subseteq\{1,\ldots,M\}}
  \mathcal{L}_{\mathrm{val}}(S)
  +\lambda_1|S|
  +\lambda_2 C(S),
  \label{eq:ensemble-static-selection}
\end{equation}
其中$\lambda_1,\lambda_2\geq0$，$|S|$控制成员数，$C(S)$可以表示存储、单样本延迟或其他部署成本。这个组合优化通常随成员池大小指数增长，需要用排序、前向加入、后向删除或其他近似搜索。

只按单成员验证损失排序，容易保留许多行为相近的模型；按当前组合的边际改进逐个加入成员，可以同时感知成员质量与已有组合。允许重复选择某个成员时，重复次数还形成离散权重。基于模型库的贪心集成选择就是这种做法的代表：从一个初始组合出发，每轮加入使验证指标改善最大的候选，必要时使用多起点或袋化减少搜索的不稳定性\scite{caruana2004ensemble}。

成员选择与一般模型选择的区别在于，式\eqref{eq:ensemble-static-selection}固定了候选池及其预测，只搜索哪些已训练成员共同部署。若同时改变基学习算法、特征处理和超参数，搜索范围已经扩展为完整模型开发流程。候选池越大，验证集上偶然出现有利组合的机会也越多；选择所用数据不能同时充当最终测试。

分歧率、误差协方差等多样性指标可以缩小候选范围，却不能替代任务损失。一个与其他成员高度不同但经常出错的模型，可能增加分歧而降低组合质量。静态选择应直接评价选定组合的输出，并在相同成员数或资源预算下与完整集成及单模型基线比较。

\subsection{动态成员选择}
\label{sec:ensemble-dynamic-selection}

静态选择假设同一成员子集适合所有输入。若成员在输入空间的不同区域各有所长，可以先估计当前输入附近的成员能力，再决定由谁预测。\term{动态分类器选择}（dynamic classifier selection, DCS）为每个输入选择一个成员；\term{动态集成选择}（dynamic ensemble selection, DES）选择一个成员子集或给出输入相关权重。

一种基本构造是准备独立的能力估计数据。对查询输入$\bm{x}$，在该数据中寻找非空邻域$\mathcal{N}(\bm{x})$，再以成员$m$在邻域内的正确率估计
\begin{equation}
  \widehat{c}_m(\bm{x})
  =
  \frac{1}{|\mathcal{N}(\bm{x})|}
  \sum_{i\in\mathcal{N}(\bm{x})}
  \mathbb{I}\!\left\{h_m(\bm{x}_i)=y_i\right\}.
  \label{eq:ensemble-local-competence}
\end{equation}
DCS可以选择$\widehat{c}_m(\bm{x})$最大的成员，DES则保留超过阈值的成员或按能力加权。局部正确率只是能力估计的一种方式，还可以使用输出轮廓、分类间隔或由元模型学习的能力特征\scite{woods1997localaccuracy}。

这个流程新增了两个误差来源。邻域依赖输入表示与距离，高维或分布稀疏时，附近样本未必提供稳定的局部证据；能力估计器也可能在有限数据上过拟合。成员训练数据、能力估计数据与最终测试数据需要分开，所有标准化和邻域参数都在相应开发流程内确定。

动态选择不一定降低端到端延迟。若系统必须先运行全部成员才能构造输出轮廓，再选择其中几个聚合，成员计算已经发生；基于输入邻域的方法还增加了检索成本。只有能够在运行昂贵成员之前完成路由，或动态组合确实改善资源分配时，成员数减少才可能转化为实际延迟收益。

\subsection{从集成模型到紧凑模型}
\label{sec:ensemble-compression}

静态与动态选择仍依赖原成员池。另一条路线是训练一个新的\term{紧凑模型}（compact model），使其逼近完整集成的输入输出映射。早期模型压缩工作已经用大型集成的输出训练较小模型；知识蒸馏随后把这种教师到学生的传递扩展为常用训练形式\scite{bucilua2006compression,hinton2015distilling}。

设完整集成为教师$H$，紧凑模型为学生$s_{\bm{\theta}}$。对回归，可以组合真实响应损失与教师拟合损失：
\begin{equation}
  \min_{\bm{\theta}}
  \sum_{i=1}^{n}
  \left[
    \alpha\ell\!\left(y_i,s_{\bm{\theta}}(\bm{x}_i)\right)
    +(1-\alpha)
    \ell\!\left(H(\bm{x}_i),s_{\bm{\theta}}(\bm{x}_i)\right)
  \right],
  \label{eq:ensemble-compression-objective}
\end{equation}
其中$\alpha\in[0,1]$控制真实标签与教师输出的作用。分类时可以拟合教师的完整类别分布；温度缩放会使非最大类别之间的相对关系更容易进入训练信号，但温度、损失权重与部署概率仍须分别选择和校准。

学生训练输入可以来自原训练集、额外未标注样本或经过说明的合成样本。教师只在这些输入上提供目标，输入覆盖不足时，学生即使在压缩数据上高度一致，也可能在其他区域偏离教师。学生容量过小会产生逼近误差，容量很大则可能无法达到预期资源节省。

紧凑化至少需要四类证据。预测一致性衡量学生是否复现教师输出；任务性能检查学生面对真实标签的风险；概率评价检查蒸馏后的得分是否仍可解释；资源测量确认模型大小、内存和延迟是否实际下降。教师与学生高度一致不自动保证二者具有相同新样本风险，学生也可能因真实标签项、正则化或额外数据而在部分样本上优于教师。原集成的成员级解释、袋外诊断和理论条件不能直接转移到新模型。

\section{模型选择与评价}
\label{sec:ensemble-evaluation}

\subsection{方法比较}
\label{sec:ensemble-method-comparison}

表\ref{tab:ensemble-method-navigation}按成员来源、依赖和组合规则比较本章方法。相同的“100个成员”在不同方法中不是同一计算单位：Bagging的一棵深树负责完整预测，Boosting的一棵浅树只提供增量，Stacking的一个成员甚至可能是已经包含许多树的随机森林。容量与成本必须沿完整训练和预测路径衡量。

\begin{table*}[htbp]
  \centering
  \small
  \begin{tabularx}{\linewidth}{@{}>{\raggedright\arraybackslash}p{23mm}*{6}{>{\raggedright\arraybackslash}X}@{}}
    \toprule
    方法 & 成员来源 & 成员间依赖 & 组合规则 & 训练并行性 & 推理路径 & 主要失效来源 \\
    \midrule
    Bagging & 重采样或随机权重下的同类成员 & 给定原训练集后可独立生成 & 平均、投票或固定加权 & 成员间可并行 & 通常运行全部成员 & 成员稳定或共同偏差占主导 \\
    随机森林／Extra Trees & 样本、特征与阈值的随机树 & 树之间无序贯依赖 & 平均或投票 & 树间可并行 & 运行全部或选定树 & 随机化削弱单树，或相关性仍高 \\
    AdaBoost & 当前样本权重下的弱分类器 & 第$t$轮依赖前轮权重 & 训练中确定的加权投票 & 轮次串行 & 累加全部或截断成员 & 噪声记录持续获得大权重 \\
    梯度提升树 & 当前损失梯度下的回归树 & 第$t$轮依赖$F_{t-1}$ & 收缩后的加法得分 & 轮次串行，单轮可并行 & 顺序累加树输出 & 损失、轮数、树复杂度与采样失配 \\
    Stacking & 同质或异质候选模型 & 基成员可独立，元层依赖折外预测 & 学习得到的元模型 & 折与成员可部分并行 & 运行基成员后再运行元模型 & 元特征泄漏、成员冗余或元层过拟合 \\
    动态选择／紧凑化 & 已训练成员池或教师集成 & 依赖能力估计或教师输出 & 输入相关选择或新学生模型 & 依具体训练流程而定 & 选择部分成员，或只运行学生 & 路由失准、覆盖不足或逼近误差 \\
    \bottomrule
  \end{tabularx}
  \caption{集成方法的属性导航。表中成本描述方法结构，不构成固定速度排序；实际结果还受成员模型、实现、硬件、数据规模与调参预算影响。}
  \label{tab:ensemble-method-navigation}
\end{table*}

方法比较还要固定组合输出。回归的均值、分位数和分布预测不能共用一个误差口径；分类的硬标签、排序得分和概率也需要不同证据。一个集成若只改善准确率，不能据此声称概率校准或尾部代价同步改善。具体指标沿用表\ref{tab:regression-evaluation-protocol}和表\ref{tab:classification-evaluation-protocol}，本节补充集成特有的数据隔离与资源记录。

\subsection{数据隔离与泄漏控制}
\label{sec:ensemble-data-isolation}

集成流程会产生多种“没有直接用于某次拟合的预测”，但它们的职责不同：
\begin{itemize}
  \item \term{袋外预测}来自未抽中当前样本的自助成员，用于Bagging内部诊断、置换重要性或相关构造。
  \item \term{折外预测}来自未使用当前折拟合的模型，用于训练Stacking元学习器，或为其他二阶段模型提供样本外特征。
  \item \term{能力估计预测}来自动态选择的独立开发数据，用于判断某个成员在当前输入附近是否可靠。
  \item \term{验证预测}用于选择成员数、树复杂度、学习率、组合约束、阈值与早停轮数。
  \item \term{测试预测}只在所有训练、选择、校准和部署规则冻结后产生，用于评价最终流程。
\end{itemize}
图\ref{fig:ensemble-data-isolation}把五类预测放回各自的数据路径。前四类仍处在模型开发阶段，允许影响规定范围内的诊断、二阶段训练或选择；测试预测位于流程冻结之后，只负责评价最终链路。区分这些职责，比笼统地把它们都称为“样本外预测”更能暴露泄漏。

\begin{figure*}[htbp]
  \centering
  % Data roles in ensemble development and final evaluation.
% Schematic only; no dataset sizes or empirical quantities are encoded.
% Canvas: 164 mm x 82 mm. Dependencies: project palette and TikZ.
\begin{tikzpicture}[
  x=1mm,y=1mm,
  font=\footnotesize,
  text=BookInk,
  source/.style={
    rectangle,
    draw=FigureInput,
    fill=FigureInput!8,
    line width=0.7pt,
    minimum width=43mm,
    minimum height=9mm,
    text width=39mm,
    align=center,
    inner sep=1mm
  },
  condition/.style={
    rectangle,
    draw=FigureLoss,
    fill=FigureLoss!9,
    line width=0.7pt,
    minimum width=52mm,
    minimum height=9mm,
    text width=48mm,
    align=center,
    inner sep=1mm
  },
  purpose/.style={
    rectangle,
    draw=FigureOutput,
    fill=FigureOutput!6,
    line width=0.8pt,
    minimum width=45mm,
    minimum height=9mm,
    text width=41mm,
    align=center,
    inner sep=1mm
  },
  final/.style={
    source,
    line width=0.9pt
  },
  final purpose/.style={
    purpose,
    line width=0.9pt
  },
  flow/.style={
    -{Latex[length=2.2mm,width=1.5mm]},
    draw=BookMuted,
    line width=0.8pt
  }
]
  \path[use as bounding box] (0,0) rectangle (164,82);

  \node[font=\heiti\small,text=FigureInput] at (25,78) {预测或数据来源};
  \node[font=\heiti\small,text=FigureLoss] at (82,78) {保持隔离的条件};
  \node[font=\heiti\small,text=FigureOutput] at (139,78) {在开发流程中的用途};

  \node[source] (oob-source) at (25,65) {袋外预测};
  \node[condition] (oob-condition) at (82,65)
    {成员$m$未使用样本$i$训练};
  \node[purpose] (oob-purpose) at (139,65)
    {Bagging内部诊断与构造};

  \node[source] (oof-source) at (25,51) {折外预测};
  \node[condition] (oof-condition) at (82,51)
    {折内模型未使用当前留出折};
  \node[purpose] (oof-purpose) at (139,51)
    {训练Stacking元学习器};

  \node[source] (competence-source) at (25,37) {能力估计预测};
  \node[condition] (competence-condition) at (82,37)
    {能力数据不参与成员训练};
  \node[purpose] (competence-purpose) at (139,37)
    {估计局部能力并路由};

  \node[source] (validation-source) at (25,23) {验证预测};
  \node[condition] (validation-condition) at (82,23)
    {与相应候选的拟合路径隔离};
  \node[purpose] (validation-purpose) at (139,23)
    {调参、筛选、校准与早停};

  \draw[BookRule,line width=0.8pt] (1,14)--(163,14);
  \node[fill=white,inner sep=1mm,font=\heiti\small,text=BookInk] at (82,14)
    {冻结全部训练、选择、校准与部署规则};

  \node[final] (test-source) at (25,6) {测试预测};
  \node[condition,line width=0.9pt]
    (test-condition) at (82,6)
    {测试样本未进入任何开发决策};
  \node[final purpose] (test-purpose) at (139,6)
    {评价最终完整流程};

  \foreach \prefix in {oob,oof,competence,validation,test} {
    \draw[flow] (\prefix-source)--(\prefix-condition);
    \draw[flow] (\prefix-condition)--(\prefix-purpose);
  }
\end{tikzpicture}

  \caption{集成学习中的数据隔离与预测职责。每一行从预测或数据来源出发，经由相应隔离条件到达允许的用途；水平分隔线表示训练、调参、筛选、校准和部署规则已经冻结。测试数据位于分隔线之后，不能沿开发路径回流。}
  \label{fig:ensemble-data-isolation}
\end{figure*}

这些预测都可能在局部意义上“样本外”，却不能互相替代。折外矩阵参与元学习器训练，因而不是最终测试证据；袋外误差若反复用于选择随机森林参数，也已经参与模型开发；动态能力数据若同时训练成员，会高估局部能力。

泄漏还可能发生在模型拟合之前。缺失值填补、标准化、特征选择、目标编码和降维若利用了当前留出折，就会把信息带入折外预测。基成员调参、成员池筛选、元学习器调参和概率校准也要放在正确的数据层级。需要同时比较候选成员和组合规则时，可以在外层保留最终测试集，在内层训练折中完成全部预处理、调参及折外特征构造；时间、患者、用户或设备等依赖单位则按部署边界分组划分。

数据隔离不是要求为每一步机械切出一份新数据。样本有限时，可以通过嵌套交叉验证或交叉拟合复用开发样本，同时保持当前预测来自未用该样本拟合的路径。应记录每个预测由哪些样本训练的哪些模型产生，使元训练、选择与最终评价的角色可追溯。

\subsection{预测质量与资源成本}
\label{sec:ensemble-quality-resource}

集成结果至少要与三个基线比较：相同任务下的合理单模型、未简化的完整集成，以及满足同一资源限制的候选方法。只比较集成与其中最差成员，会夸大组合价值；只比较最终分数，又无法判断收益来自成员质量提升、错误依赖变化还是组合规则。

对Bagging，可以同时报告代表性单成员损失、成员预测相关或双错率、集成损失以及损失随成员数的变化。对Boosting，应区分训练损失、验证轨迹和最终测试结果，并记录早停轮数。对Stacking，应比较各成员的折外表现、折外预测依赖、元模型的新增验证收益和全链路测试结果。多样性指标是诊断量，不应脱离任务损失单独优化。

概率输出需要额外检查。平均若干过度自信的概率，不保证得到可靠概率；AdaBoost和梯度提升的累计得分也需要经过适当链接才能解释为概率。若使用后处理校准，校准器只能接收未用于拟合该预测的数据，并要把它作为部署链路的一部分，在最终测试集上报告对数损失、Brier分数、可靠性及任务决策损失。

资源比较应记录完整过程，而不是只测单个成员。训练侧至少包括总计算量、墙钟时间、峰值内存、并行资源和调参次数；部署侧包括序列化模型大小、加载内存、单样本延迟、批量吞吐及预处理和路由开销。Bagging的成员训练可以并行，但推理通常运行多个成员；Boosting轮次串行，树内统计可并行；Stacking既要运行基成员，又要运行元模型。静态选择、动态选择和紧凑化是否真正节省资源，只能由端到端测量回答。

\subsection{从问题条件到候选方法}
\label{sec:ensemble-method-selection}

当单模型对训练样本、随机种子或局部分裂明显敏感，并且可以承担多个成员的存储与预测成本时，Bagging是直接候选。树成员反复依赖少数强特征时，可以进一步比较随机森林；若阈值搜索成本较高且可以接受更强的随机化，再比较Extra Trees。是否采用这些方法，应由固定任务上的单成员变化、成员依赖和组合收益共同支持。

当任务损失能够提供有意义的逐轮修正，并且训练流程可以承担序贯依赖时，可以考虑Boosting。AdaBoost适合检查重加权分类的机制；一般损失与树成员则进入GBDT。XGBoost、LightGBM和CatBoost的选择不应只依据名称：需要分别判断是否重视正则化二阶树目标与稀疏路径、直方图及大规模稀疏统计、或类别特征的有序统计，并在当前硬件和搜索预算下实测。

当候选模型已经产生互补的样本外预测，而且可以承担交叉验证与多模型部署成本时，Stacking提供学习式组合。数据较少时，固定留出会压缩元训练样本，完整折外构造通常更充分；成员高度相关时，受约束线性组合比高容量元模型更容易控制。Super Learner的理论比较范围由候选库、损失和组合类共同决定，不能替代有限样本验证。

成员池超过部署预算时，再判断要保留什么对象。希望所有输入走同一条较短路径时采用静态成员选择；有可靠局部能力数据且路由开销可控时考虑动态选择；完整集成质量合格但任何成员子集都达不到预算时，可以训练紧凑模型。以上条件只负责缩小候选范围，最终选择仍来自相同数据制度、调参预算和资源口径下的新样本证据。


\Needspace{18\baselineskip}
\section{本章小结}
\label{sec:ensemble-summary}

多个模型能够互补，不是因为“数量多”本身，而是因为成员生成、错误依赖与组合规则共同改变了预测。Bagging通过样本、特征或分裂随机性构造可并行成员，再用平均或投票削弱一部分训练波动；Boosting让每轮成员依赖当前组合，沿选定损失逐步加入修正；Stacking使用折外预测训练元学习器，使组合规则本身也由数据决定。三类机制可以嵌套，但各自的数据依赖、计算路径和失效原因不同。

成员池形成后，还可以改变部署规模。静态成员选择为所有输入固定同一子集，动态成员选择针对当前输入估计局部能力，紧凑化则训练一个新模型逼近完整集成。前两者仍需保留被选成员，后一种方法改变了部署模型；三者不能用同一个“剪枝”名称掩盖。

集成的最终证据来自冻结流程后的新样本评价。袋外预测服务于Bagging内部诊断，折外预测服务于Stacking的元特征构造，能力估计数据服务于动态选择；它们承担的职责不同，也都不能无条件替代最终测试。比较方法时还要固定调参预算，并同时报告任务损失、概率质量、模型大小和端到端延迟。后续神经网络模型部分将在这些共同原则上，进一步讨论由初始化、数据顺序与训练过程产生的神经网络集成。
